\documentclass[12pt,a4paper,UTF8]{ctexart}
\usepackage[a4paper,top=22mm,bottom=22mm,left=22mm,right=22mm,headheight=15pt]{geometry}
\newcommand{\Needspace}[1]{\par\begingroup\dimen0=#1\relax\dimen2=\pagegoal\advance\dimen2 by-\pagetotal\ifdim\dimen0>\dimen2\newpage\fi\endgroup}
\newcommand{\needspace}[1]{\par\begingroup\skip0=\lastskip\vskip-\skip0\dimen0=#1\relax\advance\dimen0 by\skip0\vskip\dimen0\penalty-100\vskip-\dimen0\vskip\skip0\penalty9999\endgroup}
\usepackage{amsmath,amssymb,mathtools,bm}
\usepackage{graphicx,xcolor,booktabs,longtable,array,enumitem,fancyhdr,caption,float,microtype,setspace}
\usepackage[hidelinks,unicode]{hyperref}
\usepackage[most]{tcolorbox}
\definecolor{ContestBlue}{HTML}{17365D}
\hypersetup{pdftitle={第056套},pdfauthor={}}
\setstretch{1.24}
\setlength{\parindent}{2em}
\setlength{\parskip}{0.35em}
\setlength{\emergencystretch}{2em}
\setlist{itemsep=0.35em,topsep=0.4em}
\captionsetup{font=small,labelfont=bf,justification=centering}
\setcounter{secnumdepth}{0}
\pagestyle{fancy}
\fancyhf{}
\fancyhead[L]{\small 第056套}
\fancyfoot[C]{\thepage}
\renewcommand{\headrulewidth}{0.4pt}
\ctexset{
  section={format=\Large\bfseries\color{ContestBlue},beforeskip=1.2em,afterskip=0.7em},
  subsection={format=\large\bfseries\color{ContestBlue},beforeskip=1.0em,afterskip=0.5em},
  subsubsection={format=\normalsize\bfseries,beforeskip=0.8em,afterskip=0.35em}
}
\usepackage{calc}
\newcounter{none}
\let\originalsubsection\subsection
\renewcommand{\subsection}{\needspace{6\baselineskip}\originalsubsection}
\makeatletter
\newcommand{\tightlist}{\setlength{\itemsep}{0pt}\setlength{\parskip}{0pt}}
\newcommand{\maxwidth}{\ifdim\Gin@nat@width>\linewidth\linewidth\else\Gin@nat@width\fi}
\newcommand{\maxheight}{\ifdim\Gin@nat@height>0.72\textheight 0.72\textheight\else\Gin@nat@height\fi}
\newcommand{\pandocbounded}[1]{\begingroup\setkeys{Gin}{width=\maxwidth,height=\maxheight,keepaspectratio}#1\endgroup}
\makeatother
\begin{document}
\begin{titlepage}
  \centering
  \vspace*{58mm}
  {\Huge\bfseries\color{ContestBlue} 第056套\par}
  \vfill
  \vspace*{28mm}
\end{titlepage}
\section{第 1 题}\label{ux7b2c-1-ux9898}

混沌摆可抽象为两根具有相同线密度的均匀细杆组成的复合摆。长摆臂 \(AB\) 在
\(O\) 点与固定支架铰接，\(OB=l_1\)、\(OA=2l_1\)。短摆臂 \(CD\) 总长为
\(l_3\)，其距端点 \(D\) 为 \(l_3/3\) 的位置与长摆臂端点 \(B\) 铰接，因此
\(BD=l_3/3\)、\(BC=2l_3/3\)。两杆夹角为 \(\theta\)。

（1）为尽量减小支架处的摩擦影响，当 \(\theta=90^\circ\)
且长摆臂水平时，要求系统能够静力平衡，求 \(l_3\)；

（2）当 \(\theta=90^\circ\)
时，在地面参考系中，长摆臂和短摆臂分别以恒定角速度
\(\omega_1\)、\(\omega_2\) 顺时针转动，求 \(C\)、\(D\)
两点的速度大小和加速度大小。

\begin{figure}[H]
\centering
\includegraphics[width=0.76\linewidth,height=\textheight,keepaspectratio,alt={图1-1　复合摆两根均匀摆臂的长度、铰点与夹角}]{../figures/question-01.pdf}
\caption*{图1-1　复合摆两根均匀摆臂的长度、铰点与夹角}
\end{figure}

\subsection{一、题型定位}\label{ux4e00ux9898ux578bux5b9aux4f4d}

静力部分使用两杆重心对O的力矩。运动部分以B为参考点，把铰点平动与短杆的绝对转动叠加；两角速度都相对地面给出。

\subsection{二、知识点}\label{ux4e8cux77e5ux8bc6ux70b9}

\begin{enumerate}
\def\labelenumi{\arabic{enumi}.}
\tightlist
\item
  均匀杆质量与长度成正比，重心位于几何中点。
\item
  平面力系对固定点的力矩平衡。
\item
  刚体速度公式 \(\bm v_P=\bm v_B+\bm\omega\times\bm r_{P/B}\)。
\item
  刚体加速度公式
  \(\bm a_P=\bm a_B+\bm\alpha\times\bm r_{P/B}+\bm\omega\times(\bm\omega\times\bm r_{P/B})\)。
\item
  匀角速转动时切向加速度为零，但向心加速度仍然存在。
\end{enumerate}

\subsection{三、物理图像}\label{ux4e09ux7269ux7406ux56feux50cf}

取 \(x\) 轴水平向右、\(y\) 轴竖直向上。当 \(\theta=90^\circ\) 且 \(AB\)
水平时，令 \(OB\) 沿 \(+x\) 方向，短杆 \(CD\) 沿竖直方向，\(D\) 在 \(B\)
上方、\(C\) 在 \(B\) 下方。顺时针角速度对应

\[
\bm\omega_1=-\omega_1\bm e_z,\qquad
\bm\omega_2=-\omega_2\bm e_z. \tag{1}
\]

\begin{figure}[H]
\centering
\includegraphics[width=0.685\linewidth,height=\textheight,keepaspectratio,alt={图1-2　同一铰点的速度和加速度合成}]{../figures/solution-01.pdf}
\caption*{图1-2　同一铰点的速度和加速度合成}
\end{figure}

图 1-2 中，长杆重心位于 \(A\)、\(B\) 的中点，其横坐标为
\(-l_1/2\)；短杆重心位于 \(CD\) 中点，相对 \(B\) 向下
\(l_3/6\)，但其重力作用线仍经过横坐标
\(x=l_1\)。这两个水平力臂是第（1）问的计算依据。

\subsection{四、解题入口}\label{ux56dbux89e3ux9898ux5165ux53e3}

设两杆材料相同、线密度均为 \(\lambda\)。长杆长度为 \(3l_1\)，质量为

\[
M_1=3\lambda l_1. \tag{2}
\]

短杆质量为

\[
M_2=\lambda l_3. \tag{3}
\]

第（2）问先求铰点 \(B\)：

\[
\bm r_{B/O}=l_1\bm e_x. \tag{4}
\]

再由短杆的分点关系写出

\[
\bm r_{D/B}=\frac{l_3}{3}\bm e_y\qquad
\bm r_{C/B}=-\frac{2l_3}{3}\bm e_y. \tag{5}
\]

这一步把几何长度和方向同时固定，后续叉乘的正负号便不再依靠想象判断。

\subsection{五、解析}\label{ux4e94ux89e3ux6790}

\needspace{3\baselineskip}
\textbf{（一）静力平衡求 \(l_3\)}

长杆重心 \(G_1\) 的坐标为

\[
x_{G_1}=\frac{x_A+x_B}{2}
=\frac{-2l_1+l_1}{2}
=-\frac{l_1}{2}. \tag{6}
\]

长杆重力对 \(O\) 的力矩大小为

\[
\begin{aligned}
\tau_1
&=M_1g\left|x_{G_1}\right|\\
&=(3\lambda l_1)g\cdot\frac{l_1}{2}\\
&=\frac32\lambda g l_1^2.
\end{aligned} \tag{7}
\]

该力矩使长杆左端下沉，取逆时针为正时 \(\tau_1>0\)。

短杆竖直，短杆重心虽然在 \(B\) 下方
\(l_3/6\)，但重力作用线的水平坐标就是 \(x_B=l_1\)。所以短杆重力对 \(O\)
的力矩大小为

\[
\begin{aligned}
\tau_2
&=M_2g l_1\\
&=(\lambda l_3)g l_1.
\end{aligned} \tag{8}
\]

它使右侧下沉，方向与 \(\tau_1\) 相反。静力平衡要求

\[
\tau_1-\tau_2=0. \tag{9}
\]

逐项代入：

\[
\frac32\lambda g l_1^2-\lambda l_3g l_1=0. \tag{10}
\]

约去非零公因子 \(\lambda g l_1\)：

\[
\frac32l_1-l_3=0. \tag{11}
\]

因此

\[
\boxed{l_3=\frac32l_1}. \tag{12}
\]

\needspace{3\baselineskip}
\textbf{（二）铰点 \(B\) 的速度和加速度}

长杆匀速绕 \(O\) 顺时针转动，故

\[
\begin{aligned}
\bm v_B
&=\bm\omega_1\times\bm r_{B/O}\\
&=(-\omega_1\bm e_z)\times(l_1\bm e_x)\\
&=-\omega_1l_1\bm e_y,
\end{aligned} \tag{13}
\]

即 \(B\) 瞬时竖直向下运动。由于 \(\omega_1\) 恒定，\(B\) 只有指向 \(O\)
的向心加速度：

\[
\bm a_B=-\omega_1^2\bm r_{B/O}
=-\omega_1^2l_1\bm e_x. \tag{14}
\]

\needspace{3\baselineskip}
\textbf{（三）端点 \(D\) 的速度}

短杆相对地面的角速度为 \(-\omega_2\bm e_z\)，所以

\[
\begin{aligned}
\bm v_D
&=\bm v_B+\bm\omega_2\times\bm r_{D/B}\\
&=-\omega_1l_1\bm e_y
+(-\omega_2\bm e_z)\times\left(\frac{l_3}{3}\bm e_y\right)\\
&=\frac{\omega_2l_3}{3}\bm e_x-\omega_1l_1\bm e_y.
\end{aligned} \tag{15}
\]

两个分量互相垂直，故速度大小为

\[
\boxed{v_D=\sqrt{\omega_1^2l_1^2+\frac{\omega_2^2l_3^2}{9}}}. \tag{16}
\]

\needspace{3\baselineskip}
\textbf{（四）端点 \(C\) 的速度}

同理，

\[
\begin{aligned}
\bm v_C
&=\bm v_B+\bm\omega_2\times\bm r_{C/B}\\
&=-\omega_1l_1\bm e_y
+(-\omega_2\bm e_z)\times\left(-\frac{2l_3}{3}\bm e_y\right)\\
&=-\frac{2\omega_2l_3}{3}\bm e_x-\omega_1l_1\bm e_y.
\end{aligned} \tag{17}
\]

所以

\[
\boxed{v_C=\sqrt{\omega_1^2l_1^2+\frac{4\omega_2^2l_3^2}{9}}}. \tag{18}
\]

\needspace{3\baselineskip}
\textbf{（五）端点 \(D\) 的加速度}

短杆也作匀角速转动，\(\bm\alpha_2=\bm0\)。因此相对 \(B\)
的加速度只有向心项：

\[
\begin{aligned}
\bm a_D
&=\bm a_B-\omega_2^2\bm r_{D/B}\\
&=-\omega_1^2l_1\bm e_x
-\frac{\omega_2^2l_3}{3}\bm e_y.
\end{aligned} \tag{19}
\]

故

\[
\boxed{a_D=\sqrt{\omega_1^4l_1^2+\frac{\omega_2^4l_3^2}{9}}}. \tag{20}
\]

\needspace{3\baselineskip}
\textbf{（六）端点 \(C\) 的加速度}

\[
\begin{aligned}
\bm a_C
&=\bm a_B-\omega_2^2\bm r_{C/B}\\
&=-\omega_1^2l_1\bm e_x
+\frac{2\omega_2^2l_3}{3}\bm e_y,
\end{aligned} \tag{21}
\]

从而

\[
\boxed{a_C=\sqrt{\omega_1^4l_1^2+\frac{4\omega_2^4l_3^2}{9}}}. \tag{22}
\]

这里没有额外的科里奥利项。公式
\(\bm a_P=\bm a_B+\bm\alpha\times\bm r+\bm\omega\times(\bm\omega\times\bm r)\)
本来就是用短杆相对地面的绝对角速度描述刚体上两点的关系；只有在另取一个随动转动坐标系，并把点相对该坐标系的运动单独拆开时，才会出现科里奥利项。

量纲检查也很直接：速度式中每一项都是 \(\omega l\)，加速度式中每一项都是
\(\omega^2l\)。若令 \(\omega_2=0\)，\(C\)、\(D\) 与 \(B\)
具有相同的瞬时速度、加速度；若令 \(\omega_1=0\)，结果退化为短杆绕固定点
\(B\) 的匀速转动。这两个极限均与直观一致。

还可以用瞬心观点复核速度方向。长杆的瞬时转动中心是 \(O\)，因此
\(\bm v_B\perp OB\)；短杆上两点的速度差必须满足

\[
\bm v_D-\bm v_C
=\bm\omega_2\times\bm r_{D/C}.
\]

由式（15）、（17）相减，

\[
\begin{aligned}
\bm v_D-\bm v_C
&=\left(\frac13+\frac23\right)\omega_2l_3\bm e_x\\
&=\omega_2l_3\bm e_x,
\end{aligned}
\]

而 \(\bm r_{D/C}=l_3\bm e_y\)，代入右侧同样得到
\(\omega_2l_3\bm e_x\)。这项复核同时确认了两端速度的水平分量符号和
\(1:2\) 比例。

若题目把``恒定角速度''改成瞬时角速度，同时给出角加速度
\(\alpha_1\)、\(\alpha_2\)，则铰点还应有切向项

\[
\bm a_B=-\omega_1^2l_1\bm e_x-\alpha_1l_1\bm e_y,
\]

短杆端点还应分别加上
\(\bm\alpha_2\times\bm r_{P/B}\)。把这些项与当前结果对比，可以看清题设``匀角速''究竟删去了哪些物理量，而不是笼统地说加速度不变。

静力矩还可直接从总质心位置核验。系统总质量为
\(\lambda(3l_1+l_3)\)，总质心横坐标满足

\[
x_G=\frac{(3\lambda l_1)(-l_1/2)+(\lambda l_3)l_1}
{\lambda(3l_1+l_3)}.
\]

固定铰点只提供通过 \(O\) 的支承力，重力合力若不通过 \(O\)
就会产生净力矩。因此平衡等价于 \(x_G=0\)。令分子为零：

\[
-\frac32l_1^2+l_1l_3=0,
\]

再次得到
\(l_3=3l_1/2\)。质心法把两个力矩合并成一个几何条件，适合推广到多根均匀杆的平衡问题。

把第（1）问结果代回第（2）问，可以得到只含 \(l_1\)
的端点结果，便于数值使用：

\[
\begin{aligned}
v_D&=l_1\sqrt{\omega_1^2+\frac14\omega_2^2},\\
v_C&=l_1\sqrt{\omega_1^2+\omega_2^2},\\
a_D&=l_1\sqrt{\omega_1^4+\frac14\omega_2^4},\\
a_C&=l_1\sqrt{\omega_1^4+\omega_2^4}.
\end{aligned}
\]

这里 \(v_C\ge v_D\)、\(a_C\ge a_D\)，因为 \(C\) 离铰点 \(B\) 更远。若
\(\omega_1=\omega_2=\omega\)，则 \(v_C=\sqrt2\omega l_1\)，而
\(a_C=\sqrt2\omega^2l_1\)；速度和加速度具有相同的几何合成系数，但量纲分别是
\(\omega l\) 与 \(\omega^2l\)。

还可检查两端点相对加速度：

\[
\begin{aligned}
\bm a_D-\bm a_C
&=-\frac{\omega_2^2l_3}{3}\bm e_y
-\frac{2\omega_2^2l_3}{3}\bm e_y\\
&=-\omega_2^2l_3\bm e_y.
\end{aligned}
\]

由于 \(\bm r_{D/C}=l_3\bm e_y\)，刚体关系要求

\[
\bm a_D-\bm a_C=-\omega_2^2\bm r_{D/C},
\]

恰与上式相同。注意铰点 \(B\)
的平动加速度在相减时自动消去；因此端点差量是检验短杆自身转动项的理想工具。速度差核验角速度的一次方，加速度差核验角速度的平方，两者同时通过才说明合成公式完整。

\textbf{线密度条件。}
原题的比值答案使用两杆具有同一线密度。若只说两杆各自均匀，而没有规定线密度相同，结果应保留\(\lambda_1,\lambda_2\)，不能直接消去质量。图中O放在实际铰接点，杆的重力作用于各自中点；\(90^\circ\)标在两杆之间，不能画成杆与坐标轴的夹角。

\needspace{3\baselineskip}
\textbf{用速度瞬心给出另一种方向检查}

在所画姿态下，\(B=(l_1,0)\)，短杆竖直，点 \(D\) 位于 \(B\) 上方。\(B\)
的速度竖直向下，所以短杆速度瞬心 \(K\) 必在过 \(B\) 的水平线上。设
\(K=(x_K,0)\)，短杆绝对角速度为顺时针 \(-\omega_2\)，由
\[\boldsymbol v_B=-\omega_2\boldsymbol e_z\times[(l_1-x_K)\boldsymbol e_x]\]
和 \(\boldsymbol v_B=-\omega_1l_1\boldsymbol e_y\) 得
\[x_K=l_1\left(1-\frac{\omega_1}{\omega_2}\right).\] 于是
\[KD^2=\left(\frac{\omega_1l_1}{\omega_2}\right)^2+\left(\frac{l_3}{3}\right)^2,\]
\[KC^2=\left(\frac{\omega_1l_1}{\omega_2}\right)^2+\left(\frac{2l_3}{3}\right)^2.\]
用
\(v_D=|\omega_2|KD,v_C=|\omega_2|KC\)，重新得到正文两种速度大小。这次核对只使用一个瞬心及几何距离，没有逐项合成端点速度。

当 \(\omega_2=\omega_1\) 时，瞬心正好位于
\(O\)。在这个瞬间两杆上的速度都可看作绕 \(O\)
转动，但两杆的加速度一般仍要由各自运动方程确定；速度瞬心不一定是固定铰点，不能直接写成所有点都绕
\(K\) 做匀速圆周运动。若
\(\omega_2=0\)，短杆作瞬时平移，瞬心在无穷远，应回到
\(\boldsymbol v_C=\boldsymbol v_D=\boldsymbol v_B\)
的表达，而不是给上述除法代入零。

\subsection{六、关键洞察}\label{ux516dux5173ux952eux6d1eux5bdf}

静力部分中，短杆重心相对 \(B\) 的竖直偏移不改变其重力对 \(O\)
的力臂；力臂是重力作用线到 \(O\) 的水平距离。运动学部分中，\(\omega_2\)
已经是短杆对地绝对角速度，不能再把 \(\omega_1\) 加到 \(\omega_2\) 上。

\subsection{七、易错点}\label{ux4e03ux6613ux9519ux70b9}

\begin{enumerate}
\def\labelenumi{\arabic{enumi}.}
\tightlist
\item
  把长杆重心误认为在 \(O\) 点；实际上 \(AB\) 不以 \(O\) 为中点。
\item
  把 \(B\) 当成短杆中点，漏掉 \(1:2\) 的分点关系。
\item
  只算短杆绕 \(B\) 的相对速度，漏掉 \(B\) 的平动速度。
\item
  匀角速被误解为``加速度为零''。
\item
  求题目要求的大小时直接把带符号分量相加，而没有先作平方和。
\end{enumerate}

\subsection{八、课堂讲解建议}\label{ux516bux8bfeux5802ux8bb2ux89e3ux5efaux8bae}

先让学生只在图上标出三个位置矢量
\(\bm r_{B/O}\)、\(\bm r_{D/B}\)、\(\bm r_{C/B}\)，再进行叉乘。每写完一个叉乘，都用顺时针瞬时运动方向做一次几何核对。静力平衡可让学生先用``左边重力矩''和``右边重力矩''口算方向，再进入代数。

\subsection{九、变式训练}\label{ux4e5dux53d8ux5f0fux8badux7ec3}

若短杆与长杆的线密度分别为 \(\lambda_2\)、\(\lambda_1\)，则平衡条件改为

\[
\frac32\lambda_1gl_1^2=\lambda_2gl_1l_3,
\]

即 \(l_3=3\lambda_1l_1/(2\lambda_2)\)。进一步可令
\(\omega_1\)、\(\omega_2\) 随时间变化，要求写出含
\(\alpha_1\)、\(\alpha_2\) 的端点加速度分量，以检验对切向加速度的掌握。

\needspace{3\baselineskip}
\textbf{含角加速度时的端点分量}

仍以顺时针的\(\omega_1,\omega_2\)为正，角加速度\(\alpha_i=\dot\omega_i\)可带符号。铰点加速度为
\[\boldsymbol a_B=(-\omega_1^2l_1,-\alpha_1l_1).\]
对D、C分别加入短杆的切向项后： \[\boldsymbol a_D=
\left(-\omega_1^2l_1+\frac{\alpha_2l_3}{3},
-\alpha_1l_1-\frac{\omega_2^2l_3}{3}\right),\] \[\boldsymbol a_C=
\left(-\omega_1^2l_1-\frac{2\alpha_2l_3}{3},
-\alpha_1l_1+\frac{2\omega_2^2l_3}{3}\right).\]
加速度大小须用这两个分量作平方和，切向项一般与原有分量不再正交。令\(\alpha_1=\alpha_2=0\)恢复原答案。

若短杆保持对地方向不变，则\(\omega_2=0\)，它是平动，而不是绕B静止。两端都有\(\boldsymbol v_B,\boldsymbol a_B\)。用这一个条件即可区分``对地角速度''与``相对长杆角速度''的不同题设。

\needspace{3\baselineskip}
\textbf{用有向角度写端点速度和加速度}

在地面系取杆1有向角为 \(\theta_1\)、杆2有向角为
\(\theta_2\)，每根杆单位向量为
\(\boldsymbol t_i=(\cos\theta_i,\sin\theta_i)\)。令
\(\boldsymbol n_i=(-\sin\theta_i,\cos\theta_i)\)，并记 \(l_2=BC\)，则
\[\boldsymbol r_B=l_1\boldsymbol t_1,\qquad
\boldsymbol r_C=l_1\boldsymbol t_1+l_2\boldsymbol t_2.\] 逐项求导得到
\[\boldsymbol v_C=l_1\omega_1\boldsymbol n_1+l_2\omega_2\boldsymbol n_2,\]
\[\boldsymbol a_C=l_1(\alpha_1\boldsymbol n_1-\omega_1^2\boldsymbol t_1)
+l_2(\alpha_2\boldsymbol n_2-\omega_2^2\boldsymbol t_2).\]
若题图两杆互相垂直，那么 \(\boldsymbol n_1,\boldsymbol n_2\)
也互相垂直，故速度大小满足 \[v_C^2=l_1^2\omega_1^2+l_2^2\omega_2^2.\]
加速度中每根杆的切向项和法向项正交，但来自不同杆的项一般并不正交；不能把四项平方直接相加。先投影到同一坐标系再合成。

例如取瞬时 \(\boldsymbol t_1=(1,0),\boldsymbol t_2=(0,1)\)，则
\[\boldsymbol v_C=(-l_2\omega_2,l_1\omega_1),\]
\[\boldsymbol a_C=(-l_1\omega_1^2-l_2\alpha_2,\,
l_1\alpha_1-l_2\omega_2^2).\] 若改用第二杆相对第一杆的角速度
\(\omega_{\rm rel}\)，地面角速度必须写成
\(\omega_2=\omega_1+\omega_{\rm rel}\)。这只是坐标定义的变化，不能把绝对角速度换成相对值后仍使用上面的公式而不补运输项。

\needspace{3\baselineskip}
\textbf{静力平衡是否稳定}

第（1）问只求出了净力矩为零的条件，不能由此断言长杆会自动回到水平。现在允许两铰点自由转动，不再施加第（2）问的给定角速度。取长杆
\(OB\) 与水平线的有向角为 \(\psi\)，短杆的重心方向相对竖直向下偏转
\(\phi\)。两杆重心高度为 \[y_{G_1}=-\frac{l_1}{2}\sin\psi,\qquad
y_{G_2}=l_1\sin\psi-\frac{l_3}{6}\cos\phi.\] 因此系统重力势能为
\[U=g\left(-\frac32\lambda_1l_1^2+\lambda_2l_1l_3\right)\sin\psi
-\frac{\lambda_2gl_3^2}{6}\cos\phi.\] 代入平衡长度
\(l_3=3\lambda_1l_1/(2\lambda_2)\)，第一项完全消失，留下
\[U=-\frac{M_2gl_3}{6}\cos\phi,\qquad M_2=\lambda_2l_3.\]
短杆重心在B下方时势能最低，微小偏转具有回复力矩；长杆角度 \(\psi\)
不进入势能，属于中性平衡。短杆保持竖直时，长杆在不同角度均可静止，而不只在水平位置平衡。这个结论要求两铰点理想、杆之间没有碰撞限制。

在长杆水平、短杆竖直向下的姿态附近，两杆的线性小振动也可分开检验。长杆绕O的转动惯量为
\[I_{1,O}=\lambda_1\int_{-2l_1}^{l_1}x^2\,dx=3\lambda_1l_1^3.\]
短杆绕B的转动惯量用平行轴定理得到
\[I_{2,B}=M_2\left(\frac{l_3^2}{12}+\frac{l_3^2}{36}\right)
=\frac{M_2l_3^2}{9}.\]
在该姿态，B的速度竖直，短杆重心相对B的微小摆动速度水平，动能交叉项在二次近似中为零。于是
\[T^{(2)}=\frac12(I_{1,O}+M_2l_1^2)\dot\psi^2
+\frac12I_{2,B}\dot\phi^2,\qquad
U^{(2)}=\frac{M_2gl_3}{12}\phi^2+\text{常量}.\] 线性方程给出
\[\ddot\psi=0,\qquad
\ddot\phi+\frac{3g}{2l_3}\phi=0.\] 相同线密度时
\(l_3=3l_1/2\)，短杆小振动角频率为
\(\sqrt{g/l_1}\)，长杆对应零频率模态。若给长杆一个微小初角速度，它不会由重力恢复到原角度；``有平衡姿态''和``有回复振动''须分别判断。

\newpage

\section{第 2 题}\label{ux7b2c-2-ux9898}

一定量理想气体经历如下 Brayton 循环：\(1\to2\) 为绝热压缩，\(2\to3\)
为等压膨胀，\(3\to4\) 为绝热膨胀，\(4\to1\) 为等压压缩。气体物质的量为
\(\nu\)，定容摩尔热容为 \(C_V=5R/2\)。求：

（1）循环中气体吸收的热量 \(Q_1\)；

（2）循环中气体放出的热量 \(Q_2\)；

（3）循环效率 \(\eta\)；

（4）\(T_1\)、\(T_2\) 与压强比 \(P_2/P_1\) 的关系；

（5）把效率写成 \(T_1\)、\(T_2\) 或压强比的函数。

\begin{figure}[H]
\centering
\includegraphics[width=0.601\linewidth,height=\textheight,keepaspectratio,alt={图2-1　Brayton 循环在 P-V 图上的四个过程}]{../figures/question-02.pdf}
\caption*{图2-1　Brayton 循环在 \(P\)-\(V\) 图上的四个过程}
\end{figure}

\subsection{一、题型定位}\label{ux4e00ux9898ux578bux5b9aux4f4d-1}

两段可逆绝热过程建立温度比，两段等压过程提供吸热和放热。由同一个压比消去温度，才能得到循环效率。

\subsection{二、知识点}\label{ux4e8cux77e5ux8bc6ux70b9-1}

\begin{enumerate}
\def\labelenumi{\arabic{enumi}.}
\tightlist
\item
  \(C_P=C_V+R\) 与热容比 \(\gamma=C_P/C_V\)。
\item
  等压过程热量 \(Q=\nu C_P\Delta T\)。
\item
  可逆绝热关系
  \(TV^{\gamma-1}=\text{常量}\)、\(TP^{(1-\gamma)/\gamma}=\text{常量}\)。
\item
  热机效率 \(\eta=1-Q_{\rm out}/Q_{\rm in}\)。
\item
  同一高压线、低压线以及两条绝热线的比例关系。
\end{enumerate}

\subsection{三、物理图像}\label{ux4e09ux7269ux7406ux56feux50cf-1}

\(1\to2\) 和 \(3\to4\) 没有热量交换；\(2\to3\) 在高压 \(P_2=P_3\)
下吸热；\(4\to1\) 在低压 \(P_4=P_1\) 下放热。用 \(T\)-\(S\)
图表示时，两条绝热线是竖直线，因而两个等压过程的温差具有相同比例。

\begin{figure}[H]
\centering
\includegraphics[width=0.548\linewidth,height=\textheight,keepaspectratio,alt={图2-2　S为横轴的Brayton循环}]{../figures/solution-02.pdf}
\caption*{图2-2　S为横轴的Brayton循环}
\end{figure}

\subsection{四、解题入口}\label{ux56dbux89e3ux9898ux5165ux53e3-1}

由 \(C_V=5R/2\)，先得到

\[
C_P=C_V+R=\frac72R, \tag{1}
\]

以及

\[
\gamma=\frac{C_P}{C_V}=\frac{7/2}{5/2}=\frac75. \tag{2}
\]

绝热过程写成温度与压强的关系：

\[
\frac{T_{\rm f}}{T_{\rm i}}
=\left(\frac{P_{\rm f}}{P_{\rm i}}\right)^{(\gamma-1)/\gamma}. \tag{3}
\]

其中

\[
\frac{\gamma-1}{\gamma}
=\frac{\frac75-1}{\frac75}
=\frac{2/5}{7/5}
=\frac27. \tag{4}
\]

\subsection{五、解析}\label{ux4e94ux89e3ux6790-1}

\needspace{3\baselineskip}
\textbf{（一）吸热量}

\(2\to3\) 是等压升温过程，故

\[
\begin{aligned}
Q_1
&=\nu C_P(T_3-T_2)\\
&=\nu\cdot\frac72R(T_3-T_2).
\end{aligned} \tag{5}
\]

因此

\[
\boxed{Q_1=\frac72\nu R(T_3-T_2)}. \tag{6}
\]

\needspace{3\baselineskip}
\textbf{（二）放热量}

\(4\to1\) 是等压降温过程。若 \(Q_2\) 表示放出热量的正值，则

\[
\begin{aligned}
Q_2
&=\nu C_P(T_4-T_1)\\
&=\frac72\nu R(T_4-T_1).
\end{aligned} \tag{7}
\]

所以

\[
\boxed{Q_2=\frac72\nu R(T_4-T_1)}. \tag{8}
\]

若按``气体吸热为正''的符号约定，\(4\to1\) 的过程热量应写成
\(Q_{41}=\nu C_P(T_1-T_4)<0\)；这里的 \(Q_2=-Q_{41}>0\)。

\needspace{3\baselineskip}
\textbf{（三）效率的第一种表达式}

循环一周内能改变量为零，因此净功为

\[
W_{\rm net}=Q_1-Q_2. \tag{9}
\]

效率为

\[
\begin{aligned}
\eta
&=\frac{W_{\rm net}}{Q_1}\\
&=\frac{Q_1-Q_2}{Q_1}\\
&=1-\frac{Q_2}{Q_1}.
\end{aligned} \tag{10}
\]

代入式（6）、（8），公共因子 \(7\nu R/2\) 消去：

\[
\boxed{\eta=1-\frac{T_4-T_1}{T_3-T_2}}. \tag{11}
\]

\needspace{3\baselineskip}
\textbf{（四）绝热压缩的温压关系}

对 \(1\to2\) 使用式（3）：

\[
\frac{T_2}{T_1}
=\left(\frac{P_2}{P_1}\right)^{2/7}. \tag{12}
\]

于是

\[
\boxed{T_2=T_1\left(\frac{P_2}{P_1}\right)^{2/7}}. \tag{13}
\]

\needspace{3\baselineskip}
\textbf{（五）用两条绝热线消去 \(T_3\)、\(T_4\)}

因为 \(P_2=P_3\)、\(P_1=P_4\)，对 \(3\to4\) 有

\[
\frac{T_3}{T_4}
=\left(\frac{P_3}{P_4}\right)^{2/7}
=\left(\frac{P_2}{P_1}\right)^{2/7}. \tag{14}
\]

与式（12）比较：

\[
\frac{T_3}{T_4}=\frac{T_2}{T_1}. \tag{15}
\]

交叉相乘得

\[
T_1T_3=T_2T_4. \tag{16}
\]

令

\[
r_T=\frac{T_2}{T_1}=\frac{T_3}{T_4}, \tag{17}
\]

则 \(T_2=r_TT_1\)、\(T_3=r_TT_4\)，于是

\[
\begin{aligned}
T_3-T_2
&=r_TT_4-r_TT_1\\
&=r_T(T_4-T_1).
\end{aligned} \tag{18}
\]

因此

\[
\frac{T_4-T_1}{T_3-T_2}=\frac1{r_T}=\frac{T_1}{T_2}. \tag{19}
\]

代回效率式：

\[
\boxed{\eta=1-\frac{T_1}{T_2}}. \tag{20}
\]

再利用式（12）：

\[
\boxed{\eta=1-\left(\frac{P_1}{P_2}\right)^{2/7}}. \tag{21}
\]

量纲上，效率必须无量纲，式（20）、（21）均满足。边界上，当
\(P_2/P_1\to1\) 时，循环在 \(P\)-\(V\)
图中收缩，\(\eta\to0\)；当压强比增大时，理想模型给出效率提高。这里并不意味着实际装置可以无限提高压强比，因为压气机、涡轮材料和最高温度会引入额外限制。

也可从绝热不变量 \(TP^{(1-\gamma)/\gamma}=\text{常量}\)
直接得到式（12）、（14）。务必保持指数的符号：写成
\(T/P^{(\gamma-1)/\gamma}=\text{常量}\) 与写成
\(TP^{(1-\gamma)/\gamma}=\text{常量}\) 完全等价。

可再用每个过程的功和内能变化核对热量。对任一等压过程，

\[
W=P\Delta V=\nu R\Delta T,
\]

而

\[
\Delta U=\nu C_V\Delta T.
\]

由热力学第一定律 \(Q=\Delta U+W\)，

\[
Q=\nu(C_V+R)\Delta T=\nu C_P\Delta T,
\]

从而验证式（5）、（7）的热容选择。对两条绝热线则
\(Q=0\)，所以压缩段外界对气体做功使内能增加，膨胀段气体对外做功使内能降低；这与
\(T_2>T_1\)、\(T_3>T_4\) 一致。

净功还可以由四段功相加得到。由于循环内能总改变量为零，

\[
\sum W=\sum Q=Q_1-Q_2.
\]

在 \(P\)-\(V\) 图上循环为顺时针，包围面积为正，因此必须有
\(Q_1>Q_2\)。这给出 \(T_3-T_2>T_4-T_1\)，也与式（18）中 \(r_T>1\)
完全相符。

绝热温压关系可以从基本方程逐步推出。可逆绝热过程满足

\[
\nu C_V\,\mathrm dT+P\,\mathrm dV=0.
\]

由理想气体状态方程 \(P=\nu RT/V\)，

\[
C_V\frac{\mathrm dT}{T}+R\frac{\mathrm dV}{V}=0.
\]

积分后得到

\[
\ln T+(\gamma-1)\ln V=\text{常量},
\]

即 \(TV^{\gamma-1}=\text{常量}\)。再利用 \(V=\nu RT/P\)：

\[
T\left(\frac{T}{P}\right)^{\gamma-1}=\text{常量}.
\]

所以 \(T^\gamma/P^{\gamma-1}=\text{常量}\)，最终得到

\[
\frac{T_2}{T_1}=\left(\frac{P_2}{P_1}\right)^{(\gamma-1)/\gamma}.
\]

这条推导能防止把 \(TV^{\gamma-1}\) 与 \(TP^{(1-\gamma)/\gamma}\)
的指数混写。

取一个具体压强比作数值检查。若 \(P_2/P_1=2^7=128\)，则

\[
\frac{T_2}{T_1}=128^{2/7}=4,
\]

效率为

\[
\eta=1-\frac14=\frac34.
\]

再任选 \(T_1=300\,\mathrm K\)、\(T_4=500\,\mathrm K\)，则
\(T_2=1200\,\mathrm K\)、\(T_3=2000\,\mathrm K\)。两个等压温差分别为
\(800\,\mathrm K\) 和 \(200\,\mathrm K\)，所以

\[
1-\frac{Q_2}{Q_1}=1-\frac{200}{800}=\frac34.
\]

温度比例法与热量定义法给出同一数值，说明式（18）中的公共比例因子使用正确。

\textbf{循环净功和适用条件。}
令压比\(\pi=P_2/P_1>1\)，\(\chi=\pi^{2/7}\)。等熵段有\(T_2=\chi T_1\)、\(T_4=T_3/\chi\)，因此
\[W=\frac{7R}{2}[(T_3-T_2)-(T_4-T_1)]
=\frac{7R}{2}(T_3-\chi T_1)(1-\chi^{-1}).\]
正向热机需\(T_3>T_2\)。图的TS坐标以S为横轴、T为纵轴，两条绝热可逆段是竖直线；等压段满足\(\mathrm dS=C_p\mathrm dT/T\)，应画成指数曲线。\(C_p=7R/2\)假定双原子气体的振动自由度尚未激发，不能不加说明地推广到任意高温。

\subsection{六、关键洞察}\label{ux516dux5173ux952eux6d1eux5bdf-1}

两条绝热线跨越的是同一对高、低压，因此温度比相同。效率化简的计算时须先区分分别求四个温度，而是识别

\[
\frac{T_2}{T_1}=\frac{T_3}{T_4}.
\]

这个比例把两个等压过程的温差直接联系起来。

\subsection{七、易错点}\label{ux4e03ux6613ux9519ux70b9-1}

\begin{enumerate}
\def\labelenumi{\arabic{enumi}.}
\tightlist
\item
  把等压热量写成 \(\nu C_V\Delta T\)。
\item
  把放热量 \(Q_2\) 的``大小''和过程热量的负号混为一谈。
\item
  将绝热指数误写成 \((\gamma-1)\)，漏掉除以 \(\gamma\)。
\item
  只写 \(\eta=1-Q_2/Q_1\)，没有完成温度比和压强比化简。
\item
  未检查 \(P_2>P_1\) 时效率应满足 \(0<\eta<1\)。
\end{enumerate}

\subsection{八、课堂讲解建议}\label{ux516bux8bfeux5802ux8bb2ux89e3ux5efaux8bae-1}

先让学生在 \(P\)-\(V\)
图上用两种颜色标出``绝热''和``等压''，并分别回答``哪里吸热、哪里放热''。随后只给出式（15），要求学生自行把
\(T_3-T_2\) 提取成
\(r_T(T_4-T_1)\)。这样能把计算的计算依据从公式记忆转向比例结构。

\subsection{九、变式训练}\label{ux4e5dux53d8ux5f0fux8badux7ec3-1}

若气体为一般定热容理想气体，保留 \(\gamma\) 不代数值，可得

\[
\eta=1-\left(\frac{P_1}{P_2}\right)^{(\gamma-1)/\gamma}.
\]

可进一步给定压强比和最低温度，要求反求压缩终温；或给定效率，要求反求所需压强比，并讨论结果是否符合实际温度上限。

\needspace{3\baselineskip}
\textbf{固定最高、最低温度时的最大净功}

保持可逆Brayton循环的\(T_1,T_3\)固定，设\(\chi=(P_2/P_1)^{(\gamma-1)/\gamma}\)，则
\[T_2=\chi T_1,\qquad T_4=T_3/\chi.\] 每摩尔净功为
\[w=C_p[T_3(1-\chi^{-1})-T_1(\chi-1)].\]
可行热机要求\(1<\chi<T_3/T_1\)。求导得
\[\frac{\mathrm dw}{\mathrm d\chi}
=C_p(T_3/\chi^2-T_1).\] 最大功对应
\[\chi_{\rm opt}=\sqrt{T_3/T_1},\qquad
\left(\frac{P_2}{P_1}\right)_{\rm opt}
=\left(\frac{T_3}{T_1}\right)^{\gamma/[2(\gamma-1)]}.\]
对\(\gamma=7/5\)，压力比的指数为\(7/4\)，最大净功是
\[w_{\max}=C_p(\sqrt{T_3}-\sqrt{T_1})^2.\]
此时效率为\(1-\sqrt{T_1/T_3}\)，低于同温限Carnot效率。继续提高压力比会提高理想效率公式，却压缩吸热温差；到热机可行边界时净功趋零，不能把最大效率等同于最大功。

\needspace{3\baselineskip}
\textbf{给定最高、最低温度的数值循环}

设 \(T_3=4T_1\)，以 \(\chi=(P_2/P_1)^{2/7}\) 表示绝热温比。净比功为
\[w=C_pT_1[5-\chi-4/\chi].\] 可行吸热范围是 \(1<\chi<4\)。求导：
\[\frac{\mathrm dw}{\mathrm d\chi}=C_pT_1(-1+4/\chi^2),\] 故最大功在
\(\chi=2\)，压比为 \(P_2/P_1=2^{7/2}=8\sqrt2\)。四状态温度为
\[T_1,\quad T_2=2T_1,\quad T_3=4T_1,\quad T_4=2T_1.\] 热量和净功为
\[Q_{\rm in}=2C_pT_1,\qquad Q_{\rm out}=C_pT_1,\qquad
w_{\max}=C_pT_1.\] 最大功时效率为 \(1/2\)，而相同温度界限的Carnot效率为
\(3/4\)。Brayton循环在不同温度吸、放热，不能把两个界限温度直接代入Carnot式。

若增大 \(\chi\) 至接近4，效率 \(1-1/\chi\) 趋近 \(3/4\)，但吸热量
\(C_pT_1(4-\chi)\)
与净功都趋于零。最大效率的退化极限和最大净功工况因此不同。课堂计算可同时给出
\(Q_{\rm in},Q_{\rm out},w\)，避免只比较一个效率数字。

\needspace{3\baselineskip}
\textbf{四状态的熵坐标怎么标在图上}

沿用 \(T_3=4T_1,\chi=2\) 的最大功例子，取状态1的熵为零。绝热段可逆，故
\[S_1=S_2=0,\qquad S_3=S_4.\] 两条等压段的熵差为
\[\Delta S_{23}=C_p\ln(T_3/T_2)=C_p\ln2,\]
\[\Delta S_{41}=C_p\ln(T_1/T_4)=-C_p\ln2.\] 因此TS图的四点可写成
\[1=(0,T_1),\quad2=(0,2T_1),\]
\[3=(C_p\ln2,4T_1),\quad4=(C_p\ln2,2T_1).\] 等压过程满足
\(T=T_{\rm start}\exp[(S-S_{\rm start})/C_p]\)，所以两条曲线应为指数曲线，不能把它们画成等温横线。上、下曲线间的面积为
\[\int_0^{C_p\ln2}(2T_1-T_1)e^{S/C_p}\,\mathrm dS=C_pT_1,\]
与循环净功一致。PV图面积与TS图面积在可逆循环中都对应净功，但它们的横轴量纲不同，不能把两幅图的斜率数值直接比较。

\newpage

\section{第 3 题}\label{ux7b2c-3-ux9898}

有 \(5\) 个彼此相邻的相同矩形木块，每块长为 \(l\)、质量为
\(m\)，置于光滑水平桌面上。质量均为 \(5m\) 的物体 \(A\)、\(B\)
由轻绳跨过桌边光滑定滑轮连接。\(A\)视为质点，起初位于最左侧木块的左端，木块与桌面间光滑，\(A\)
与木块间的动、静摩擦系数均取为
\(\mu\)，木块右侧有足够长的桌面。系统由静止释放。求：

（1）最终 \(A\) 与 \(5\) 个木块以共同加速度运动时，\(\mu\) 的范围；

（2）最终 \(A\) 与 \(4\) 个木块以共同加速度运动时，\(\mu\) 的范围；

（3）最终 \(A\) 与 \(3\) 个木块以共同加速度运动时，\(\mu\) 的范围；

（4）推广到 \(N\) 个木块，\(A\)、\(B\) 的质量均为 \(Nm\)。当
\(N\)、最终共同运动的木块数 \(K\) 都很大时，求 \(\mu\)
的近似条件。可使用 \(\sum_{i=1}^{n}1/i\simeq\ln n\)。

\begin{figure}[H]
\centering
\includegraphics[width=0.88\linewidth,height=\textheight,keepaspectratio,alt={图3-1　物体 A 依次越过相邻木块并由悬挂物体 B 牵引}]{../figures/question-03.pdf}
\caption*{图3-1　物体 \(A\) 依次越过相邻木块并由悬挂物体 \(B\) 牵引}
\end{figure}

\subsection{一、题型定位}\label{ux4e00ux9898ux578bux5b9aux4f4d-2}

A每越过一块木块，承载组的加速度便改变。是否最终锁定要累计前面各段的相对速度，不能只比较最后一段的加速度。

\subsection{二、知识点}\label{ux4e8cux77e5ux8bc6ux70b9-2}

\begin{enumerate}
\def\labelenumi{\arabic{enumi}.}
\tightlist
\item
  绳约束下 \(A\)、\(B\) 加速度大小相同。
\item
  滑动摩擦力恒为 \(\mu Nmg\)。
\item
  \(n\) 个相邻木块作为整体的加速度。
\item
  分段使用 \(v^2-v_0^2=2a\Delta x\)。
\item
  最终锁定的判据：进入第 \(K\) 段时相对速度为正，在该段长度内降为零。
\item
  调和数 \(H_n\) 及 \(H_N-H_K\simeq\ln(N/K)\)。
\end{enumerate}

\subsection{三、物理图像}\label{ux4e09ux7269ux7406ux56feux50cf-2}

当 \(A\) 与木块发生相对滑动时，摩擦力对 \(A\) 向左，对当前仍相邻的 \(n\)
个木块整体向右。\(A\)
越过最左端一块后，该块落在后方，不再把摩擦力传给右侧木块组，于是 \(n\)
依次从 \(N\) 减到 \(N-1\)、\(N-2\)\ldots\ldots。

\begin{figure}[H]
\centering
\includegraphics[width=0.755\linewidth,height=\textheight,keepaspectratio,alt={图3-2　按承载块数分段的相对速度平方}]{../figures/solution-03.pdf}
\caption*{图3-2　按承载块数分段的相对速度平方}
\end{figure}

\subsection{四、解题入口}\label{ux56dbux89e3ux9898ux5165ux53e3-2}

先在滑动阶段求 \(A\) 的恒定加速度。设绳张力为 \(T\)，取 \(A\)
向右、\(B\) 向下为正：

\[
T-\mu Nmg=Nm a_A, \tag{1}
\]

\[
Nmg-T=Nm a_A. \tag{2}
\]

两式相加：

\[
Nmg-\mu Nmg=2Nm a_A. \tag{3}
\]

所以

\[
\boxed{a_A=\frac{1-\mu}{2}g}. \tag{4}
\]

当仍有 \(n\) 个相邻木块受 \(A\) 的摩擦推动时，木块组总质量为 \(nm\)，故

\[
\boxed{a_n=\frac{\mu Nmg}{nm}=\frac{\mu N}{n}g}. \tag{5}
\]

于是该阶段 \(A\) 相对木块组的加速度为

\[
\boxed{a_{{\rm rel},n}=a_A-a_n
=\frac g2\left(1-\mu-\frac{2\mu N}{n}\right)}. \tag{6}
\]

\subsection{五、解析}\label{ux4e94ux89e3ux6790-2}

\needspace{3\baselineskip}
\textbf{（一）建立跨越若干木块后的相对速度}

记 \(u_n\) 为 \(A\) 刚进入``尚有 \(n\)
个木块''的阶段时，相对该木块组的速度。初始时系统静止，故

\[
u_N=0. \tag{7}
\]

在第 \(n\) 段中，\(A\) 要相对木块组前进一个木块长度 \(l\)
才会进入下一段。由匀变速关系，

\[
u_{n-1}^2-u_n^2=2a_{{\rm rel},n}l. \tag{8}
\]

把式（6）代入：

\[
u_{n-1}^2-u_n^2
=gl\left(1-\mu-\frac{2\mu N}{n}\right). \tag{9}
\]

从 \(n=N\) 累加到 \(n=K+1\)：

\[
\begin{aligned}
u_K^2-u_N^2
&=gl\sum_{n=K+1}^{N}
\left(1-\mu-\frac{2\mu N}{n}\right)\\
&=gl\left[(N-K)(1-\mu)
-2\mu N\sum_{n=K+1}^{N}\frac1n\right].
\end{aligned} \tag{10}
\]

利用 \(u_N=0\)：

\[
\boxed{u_K^2
=gl\left[(N-K)(1-\mu)
-2\mu N(H_N-H_K)\right]}. \tag{11}
\]

\needspace{3\baselineskip}
\textbf{（二）最终与 \(K\) 块共同运动的两个边界}

要确实进入第 \(K\) 段，进入时必须仍有正的相对速度，即

\[
u_K^2>0. \tag{12}
\]

由式（11）：

\[
(N-K)(1-\mu)>2\mu N(H_N-H_K). \tag{13}
\]

移项并提取 \(\mu\)：

\[
N-K>\mu\left[N-K+2N(H_N-H_K)\right]. \tag{14}
\]

所以右边界为

\[
\mu<\frac{N-K}{N-K+2N(H_N-H_K)}. \tag{15}
\]

进入第 \(K\) 段后，若在尚未再越过一整块之前相对速度降为零，就会与这
\(K\) 块锁定。把第 \(K\) 段也累加进去，相当于要求``若再前进完整的
\(l\)，算得的相对速度平方不为正''：

\[
u_{K-1}^2
=gl\left[(N-K+1)(1-\mu)
-2\mu N(H_N-H_{K-1})\right]\le0. \tag{16}
\]

因此

\[
(N-K+1)(1-\mu)
\le2\mu N(H_N-H_{K-1}). \tag{17}
\]

整理得左边界

\[
\mu\ge
\frac{N-K+1}
{N-K+1+2N(H_N-H_{K-1})}. \tag{18}
\]

综上，\(1\le K<N\) 时

\[
\boxed{
\frac{N-K+1}{N-K+1+2N(H_N-H_{K-1})}
\le\mu<
\frac{N-K}{N-K+2N(H_N-H_K)} }. \tag{19}
\]

等号归入左端，是因为恰在第 \(K\)
块的右边界相对速度降为零时，仍可视为最终与 \(K\)
块共同运动；右端等号对应刚进入第 \(K\) 段时相对速度已经为零，应归入
\(K+1\) 块的情形。

\needspace{3\baselineskip}
\textbf{（三）最终与全部 \(5\) 块共同运动}

此时初始便不能发生向右相对滑动。由式（6）取 \(N=n=5\)：

\[
a_{{\rm rel},5}
=\frac g2(1-\mu-2\mu)
=\frac g2(1-3\mu). \tag{20}
\]

要求 \(a_{{\rm rel},5}\le0\)：

\[
1-3\mu\le0. \tag{21}
\]

故

\[
\boxed{\mu\ge\frac13}. \tag{22}
\]

锁定后整体加速度为 \(5g/(10+5)=g/3\)，推动 \(5\) 块所需摩擦为
\(5m(g/3)=5mg/3\)，恰不超过最大静摩擦 \(\mu(5mg)\) 的条件也是
\(\mu\ge1/3\)。

\needspace{3\baselineskip}
\textbf{（四）最终与 \(4\) 块共同运动}

取 \(N=5\)、\(K=4\)。左边界用式（18）：

\[
\begin{aligned}
\mu_{\min}
&=\frac{2}{2+10(H_5-H_3)}\\
&=\frac{2}{2+10\left(\frac14+\frac15\right)}\\
&=\frac{2}{2+\frac92}\\
&=\frac{2}{\frac{13}{2}}\\
&=\frac4{13}.
\end{aligned} \tag{23}
\]

右边界用式（15）：

\[
\begin{aligned}
\mu_{\max}
&=\frac{1}{1+10(H_5-H_4)}\\
&=\frac{1}{1+10\cdot\frac15}\\
&=\frac13.
\end{aligned} \tag{24}
\]

所以

\[
\boxed{\frac4{13}\le\mu<\frac13}. \tag{25}
\]

\needspace{3\baselineskip}
\textbf{（五）最终与 \(3\) 块共同运动}

取 \(N=5\)、\(K=3\)。左边界：

\[
\begin{aligned}
\mu_{\min}
&=\frac{3}{3+10(H_5-H_2)}\\
&=\frac{3}{3+10\left(\frac13+\frac14+\frac15\right)}\\
&=\frac{3}{3+10\cdot\frac{47}{60}}\\
&=\frac{3}{3+\frac{47}{6}}\\
&=\frac{3}{\frac{65}{6}}\\
&=\frac{18}{65}.
\end{aligned} \tag{26}
\]

右边界正好是上一问左边界：

\[
\mu_{\max}=\frac4{13}. \tag{27}
\]

故

\[
\boxed{\frac{18}{65}\le\mu<\frac4{13}}. \tag{28}
\]

\needspace{3\baselineskip}
\textbf{（六）\(N\)、\(K\) 都很大时的近似式}

当 \(N\)、\(K\)
很大且相差不只一块时，式（19）的左右边界之差为高阶小量。使用

\[
H_N-H_K=\sum_{n=K+1}^{N}\frac1n
\simeq\ln\frac NK, \tag{29}
\]

得到

\[
\boxed{
\mu\simeq
\frac{N-K}{N-K+2N\ln(N/K)} }. \tag{30}
\]

若记最终保留比例 \(x=K/N\)，则

\[
\boxed{
\mu\simeq\frac{1-x}{1-x+2\ln(1/x)}}. \tag{31}
\]

为什么只比较 \(a_A\) 与 \(a_K\) 会给出错误边界？因为进入第 \(K\)
段时，\(A\) 通常已有非零相对速度。即使此时
\(a_A<a_K\)，也需要一段相对位移才能把相对速度减到零。式（16）正是把这段``刹车距离''计算在内。

锁定后的静摩擦可行性还可独立核查。\(A\)、\(B\) 与 \(K\)
块共同加速时，系统外驱动力为 \(Nmg\)，总质量为 \((2N+K)m\)，故
\(a=Ng/(2N+K)\)。推动 \(K\) 块所需摩擦与法向力之比为
\(K/(2N+K)\)。式（18）给出的动力学下界不小于这一静摩擦下界，因而最终锁定在力学上可维持。

也可以把第 \(K\) 段中的停止位置直接算出来。若
\(a_{{\rm rel},K}<0\)，从进入该段起量取相对位移 \(s\)，有

\[
0=u_K^2+2a_{{\rm rel},K}s,
\]

所以

\[
s=-\frac{u_K^2}{2a_{{\rm rel},K}}.
\]

最终与 \(K\) 块锁定要求 \(0<s\le l\)。其中 \(s>0\) 给出进入时
\(u_K^2>0\) 且相对加速度为负；\(s\le l\)
恰好等价于式（16）。这个写法能把求和不等式还原为清楚的``刹车距离''图像。

对于 \(N=5\)，各区间应首尾连接：

\[
\mu\ge\frac13\Rightarrow K=5,
\]

\[
\frac4{13}\le\mu<\frac13\Rightarrow K=4,
\]

\[
\frac{18}{65}\le\mu<\frac4{13}\Rightarrow K=3.
\]

随着 \(\mu\) 减小，\(A\)
能越过更多木块，最终共同运动的木块数单调减少；若算出的区间顺序与此相反，就说明求和上下限或端点归属有误。

大数近似还可检查两个极限。令 \(x=K/N\)。当 \(x\to1\)，设
\(x=1-\varepsilon\)，则

\[
\ln\frac1x=-\ln(1-\varepsilon)
=\varepsilon+\frac{\varepsilon^2}{2}+O(\varepsilon^3).
\]

所以

\[
\mu\simeq
\frac{\varepsilon}{\varepsilon+2\varepsilon+O(\varepsilon^2)}
\to\frac13,
\]

与``几乎全部木块留下''的临界值一致。当 \(x\to0\)
时，\(\ln(1/x)\to\infty\)，于是 \(\mu\to0\)；摩擦越弱，最终随 \(A\)
共同运动的木块比例越小。两个极限都支持式（31）的物理意义。

对第（2）问还可不调用通式，直接演示一次分段计算。初段有 \(n=5\)，

\[
a_{{\rm rel},5}=\frac g2(1-3\mu),
\]

越过第一块后

\[
u_4^2=gl(1-3\mu).
\]

进入 \(n=4\) 段后

\[
a_{{\rm rel},4}
=\frac g2\left(1-\mu-\frac{10\mu}{4}\right)
=\frac g2\left(1-\frac72\mu\right).
\]

在该段内停下要求

\[
u_4^2+2a_{{\rm rel},4}l\le0.
\]

代入并合并：

\[
gl\left[(1-3\mu)+\left(1-\frac72\mu\right)\right]
=gl\left(2-\frac{13}{2}\mu\right)\le0,
\]

得到 \(\mu\ge4/13\)；而进入第 \(4\) 段要求 \(u_4^2>0\)，即
\(\mu<1/3\)。这正好重现式（25），也展示了通式背后的两段运动。

同样地，进入 \(n=3\) 段时已经依次越过了 \(n=5\)、\(n=4\) 两段，

\[
\begin{aligned}
u_3^2
&=gl\left[\left(1-3\mu\right)
+\left(1-\frac72\mu\right)\right]\\
&=gl\left(2-\frac{13}{2}\mu\right).
\end{aligned}
\]

要求确实进入此段，\(u_3^2>0\)，给出 \(\mu<4/13\)。在 \(n=3\) 段中

\[
a_{{\rm rel},3}
=\frac g2\left(1-\mu-\frac{10\mu}{3}\right)
=\frac g2\left(1-\frac{13}{3}\mu\right).
\]

在一块长度内停下要求

\[
u_3^2+2a_{{\rm rel},3}l\le0,
\]

即

\[
gl\left(3-\frac{65}{6}\mu\right)\le0,
\]

所以
\(\mu\ge18/65\)。这样不借助通式也完整得到第（3）问区间，并清楚显示系数
\(65\) 来自前三段相对加速度的累加。

\needspace{3\baselineskip}
\textbf{相对速度平方与静摩擦分别核验}

取 \(N=5,\mu=0.29\)，每段长度为 \(l\)。以 \(w=v_{\rm rel}^2/(gl)\) 和
\(\xi=x/l\) 表示相对速度和位移。第一个阶段 \(n=5\) 的斜率为
\(0.13\)，所以 \[w(0)=0,\qquad w(1)=0.13.\] 第二个阶段 \(n=4\) 的斜率为
\(-0.015\)，在 \(\xi=2\) 时 \[w(2)=0.13-0.015=0.115>0.\]
虽然第二段相对加速度已经反向，滑动并没有立即停止，因为进入该段时已有正相对速度。第三阶段
\(n=3\) 的斜率为 \(-77/300\)，停止位移满足
\[0=0.115-\frac{77}{300}(\xi_s-2),\]
\[\xi_s=2+\frac{69}{154}=\frac{377}{154}\simeq2.448052.\]
图3-2纵轴实际放大20倍，所以三个节点的图示高度为
\(0,2.6,2.3\)；这些是绘图比例，代入运动学时仍使用 \(w=0,0.13,0.115\)。

相对滑动停止后还须检查静摩擦。此时留下 \(n=3\) 段，合运动加速度为
\(5g/13\)，维持锁定所需系数为 \(3/13\simeq0.230769\)，小于
\(0.29\)，可以锁定。若将动、静摩擦系数分别指定，就应拿 \(3/13\)
与静摩擦系数比较；不能拿第三阶段的动态减速斜率直接判断静止是否稳定。阶段能量增量、停止位置与静摩擦条件分别承担三个不同的判据。

\needspace{3\baselineskip}
\textbf{停止滑动的时间与A的实际速度}

仍用 \(\mu=0.29\) 的例子。前两段的边界相对速度为
\[u_1=\sqrt{0.13gl},\qquad u_2=\sqrt{0.115gl}.\]
在每段内相对加速度恒定，用平均相对速度乘时间求位移。第一段从静止开始，故
\[\Delta t_5=\frac{2l}{u_1}
=\frac2{\sqrt{0.13}}\sqrt{\frac lg}.\] 第二段仍跨越完整一块，其时间为
\[\Delta t_4=\frac{2l}{u_1+u_2}
=\frac2{\sqrt{0.13}+\sqrt{0.115}}\sqrt{\frac lg}.\]
第三段在块内减速到零，时间为 \[\Delta t_3=\frac{u_2}{|a_{{\rm rel},3}|}
=\frac{2\sqrt{0.115}}{77/300}\sqrt{\frac lg}.\]
三项相加是停止相对滑动的实际时间。它们使用相对速度，只用于相对位移；不能把
\(u_1,u_2\) 当成A在地面系的速度。滑动阶段A一直以 \((1-\mu)g/2=0.355g\)
加速，所以其停止滑动时的地面速度是
\[v_A=0.355g(\Delta t_5+\Delta t_4+\Delta t_3).\]
承载组此时具有同一速度，随后共同加速度变为
\(5g/13\)。速度连续，加速度可以跳变，因为摩擦从给定动摩擦值改为维持锁定所需的静摩擦值。被越过的木块失去摩擦驱动后在光滑桌面上保持各自速度，不应继续随当前承载组加速。

\subsection{六、关键洞察}\label{ux516dux5173ux952eux6d1eux5bdf-2}

本题的``题眼''是把木块数 \(n\)
当成阶段编号。每个阶段的相对加速度恒定，但各阶段不同；相对速度平方的改变量可以直接相加，最终自然出现
\(\sum1/n\)。

\subsection{七、易错点}\label{ux4e03ux6613ux9519ux70b9-2}

\begin{enumerate}
\def\labelenumi{\arabic{enumi}.}
\tightlist
\item
  把所有木块在全过程中始终视为同一整体。
\item
  只写静摩擦条件 \(\mu\ge K/(2N+K)\)，没有处理到达第 \(K\)
  段前的相对速度。
\item
  漏掉 \(A\) 的法向力为 \(Nmg\)，从而把摩擦力写成 \(\mu mg\)。
\item
  分段求和的上下限错一位，导致区间边界无法首尾相接。
\item
  不说明端点等号归属，使``最终有几块''在临界情形下含糊。
\end{enumerate}

\subsection{八、课堂讲解建议}\label{ux516bux8bfeux5802ux8bb2ux89e3ux5efaux8bae-2}

先只讨论 \(N=5\) 的第一段：让学生算出 \(a_A\)、\(a_5\)
和离开第一块时的相对速度。再追问``到了第二块，初相对速度还是零吗？''由此自然引出递推式（8）。最后再把三段式写成求和，避免直接抛出调和数公式。

\subsection{九、变式训练}\label{ux4e5dux53d8ux5f0fux8badux7ec3-2}

可把每块长度改为 \(l_n\)，则递推式成为

\[
u_{n-1}^2-u_n^2=2a_{{\rm rel},n}l_n,
\]

不再能简单提出公共长度。也可令木块质量按 \(m_n\)
变化，要求学生重新识别``受力不变、承载总质量变化''的结构。

\needspace{3\baselineskip}
\textbf{大N近似的误差与边界宽度}

保持\(x=K/N\)在\((0,1)\)的固定区间内，调和数差的展开为
\[H_N-H_K=\ln(N/K)+\frac1{2N}-\frac1{2K}+O(N^{-2}).\]
代入精确上界，若记\(D=1-x+2\ln(1/x)\)，有 \[\mu_{\rm upper}
=\frac{1-x}{D+(1-1/x)/N+O(N^{-2})}.\]
因此正文的大数式的绝对误差通常为\(O(1/N)\)，它没有分辨``恰好剩K块''和``剩K+1块''的精度。

令\(f(x)=(1-x)/D\)。两个相邻K边界相差约
\[\Delta\mu\simeq\frac{f'(x)}N,\qquad
f'(x)=\frac{2[(1-x)/x-\ln(1/x)]}{D^2}>0.\]
正号来自\(z-1>\ln z\)，其中\(z=1/x>1\)。所以保留更多木块需要更大摩擦。若K很小或\(N-K\)只有一两块，应回到精确调和数，不用固定比例展开。

\needspace{3\baselineskip}
\textbf{不等长木块的两阶段完整例子}

仍取 \(N=5\)，但最左第一块长 \(l\)，第二块长
\(2l\)，后面木块足够长。设希望最终在还剩4块时锁定，仍把 \(A\)
视作质点。第一阶段进入速度平方为 \[u_4^2=gl(1-3\mu).\]
要离开第一块，要求 \(\mu<1/3\)。在还剩4块的阶段，相对加速度为
\[a_{{\rm rel},4}=\frac g2\left(1-\frac72\mu\right).\]
先要求它能够减速，即 \(\mu>2/7\)。在第二块的长度 \(2l\) 内能降为零，要求
\[u_4^2+2a_{{\rm rel},4}(2l)\le0,\]
\[gl[(1-3\mu)+2(1-\tfrac72\mu)]\le0,\] 即
\(\mu\ge3/10\)。锁定后的加速度为 \(5g/14\)，静摩擦系数下限为 \(2/7\)，被
\(\mu\ge3/10\) 自动满足。因此范围为 \[\frac3{10}\le\mu<\frac13.\] 取
\(\mu=0.32\)，第二阶段开始的速度平方为 \(0.04gl\)，斜率为
\(-0.12g\)，停止的相对位移为 \(l/3\)，小于允许的 \(2l\)。总位移为
\(4l/3\)。
这一变式只改变长度，不能继续使用等长情况下的调和数闭式；加速度仍相同，但各段速度平方增量必须分别乘实际长度。

\newpage

\section{第 4 题}\label{ux7b2c-4-ux9898}

边长为 \(a\) 的等边三角形 \(ABC\) 的重心为
\(O\)，三个顶点分别固定电荷量均为 \(Q\)
的同号点电荷，只考虑三角形所在平面内的问题。

（1）与 \(Q\) 同号、质量为 \(m\) 的试探电荷 \(q\) 位于中线 \(AD\)
上、\(O\) 下方 \(x=\delta a\) 处。其合力写为
\(F=\alpha kQq/a^2\)，规定向上为正，求 \(\alpha\)；

（2）当 \(\delta\ll1\) 时证明 \(q\) 作小振动，并在
\(T=\beta\sqrt{ma^3/(kqQ)}\) 中求 \(\beta\)；

（3）把 \(q\) 放在底边中点 \(D\)，若
\(F_D=\alpha'kQq/a^2\)，向上为正，求 \(\alpha'\)；

（4）在整个平面内，\(q\) 共有多少个平衡位置；

（5）把 \(q\) 限制在过 \(O\) 且平行于 \(BC\) 的直线 \(PQ\) 上运动，判断
\(O\) 点平衡的稳定性；

（6）若四个等量同号电荷位于正方形四个顶点，平面内共有多少个平衡位置；

（7）若 \(2020\) 个等量同号电荷位于正 \(2020\)
边形顶点，平面内共有多少个平衡位置。

\begin{figure}[H]
\centering
\includegraphics[width=0.535\linewidth,height=\textheight,keepaspectratio,alt={图4-1　等边三角形顶点电荷与中线上的试探电荷}]{../figures/question-04.pdf}
\caption*{图4-1　等边三角形顶点电荷与中线上的试探电荷}
\end{figure}

\subsection{一、题型定位}\label{ux4e00ux9898ux578bux5b9aux4f4d-3}

先叠加三个点电荷的有向力，再展开中心附近的线性回复项。全平面的平衡点计数还需排除其他零点，不能由对称性直接写出总数。

\subsection{二、知识点}\label{ux4e8cux77e5ux8bc6ux70b9-3}

\begin{enumerate}
\def\labelenumi{\arabic{enumi}.}
\tightlist
\item
  点电荷电场的矢量叠加。
\item
  等边三角形的外接圆半径 \(a/\sqrt3\) 和重心到底边距离 \(a/(2\sqrt3)\)。
\item
  一阶泰勒展开与简谐运动判据。
\item
  势能极值和约束方向稳定性。
\item
  正多边形的旋转、镜面对称与平衡点轨道。
\end{enumerate}

\subsection{三、物理图像}\label{ux4e09ux7269ux7406ux56feux50cf-3}

建立以 \(O\) 为原点、竖直向上为 \(y\) 正方向的坐标系：

\[
A=\left(0,\frac{a}{\sqrt3}\right),\quad
B=\left(-\frac a2,-\frac{a}{2\sqrt3}\right),\quad
C=\left(\frac a2,-\frac{a}{2\sqrt3}\right). \tag{1}
\]

试探电荷在 \(O\) 下方时

\[
q=(0,-\delta a). \tag{2}
\]

由于 \(B\)、\(C\) 关于中线对称，它们的水平力分量相消，合力只沿 \(AD\)。

\begin{figure}[H]
\centering
\includegraphics[width=0.82\linewidth,height=\textheight,keepaspectratio,alt={图4-2　精确中线力和三个非中心平衡点}]{../figures/solution-04.pdf}
\caption*{图4-2　精确中线力和三个非中心平衡点}
\end{figure}

\subsection{四、解题入口}\label{ux56dbux89e3ux9898ux5165ux53e3-3}

先把所有距离除以 \(a\)，使用无量纲纵坐标 \(u=y/a=-\delta\)。\(A\) 到
\(q\) 的距离为

\[
r_A=a\left(\frac1{\sqrt3}+\delta\right). \tag{3}
\]

\(B\)、\(C\) 到 \(q\) 的竖直坐标差均为

\[
\Delta y=a\left(\frac1{2\sqrt3}-\delta\right), \tag{4}
\]

距离平方为

\[
r_B^2=r_C^2
=a^2\left[\frac14+\left(\frac1{2\sqrt3}-\delta\right)^2\right]. \tag{5}
\]

\subsection{五、解析}\label{ux4e94ux89e3ux6790-3}

\needspace{3\baselineskip}
\textbf{（一）精确求 \(\alpha(\delta)\)}

电荷 \(A\) 在 \(q\) 上方，同号排斥使 \(q\)
受到向下的力，其无量纲竖直分量为

\[
\alpha_A=-\frac{1}{\left(\frac1{\sqrt3}+\delta\right)^2}. \tag{6}
\]

来自 \(B\) 的竖直分量为

\[
F_{By}
=\frac{kQq}{r_B^2}\frac{\Delta y}{r_B}
=\frac{kQq}{a^2}
\frac{\frac1{2\sqrt3}-\delta}
{\left[\frac14+\left(\frac1{2\sqrt3}-\delta\right)^2\right]^{3/2}}. \tag{7}
\]

\(C\) 的竖直分量与之相同，所以

\[
\boxed{
\alpha(\delta)=
-\frac{1}{\left(\frac1{\sqrt3}+\delta\right)^2}
+\frac{2\left(\frac1{2\sqrt3}-\delta\right)}
{\left[\frac14+\left(\frac1{2\sqrt3}-\delta\right)^2\right]^{3/2}}
}. \tag{8}
\]

令 \(\delta=0\)，两项分别为 \(-3\) 与 \(+3\)，故
\(\alpha(0)=0\)，与中心的三重旋转对称性一致。

\needspace{3\baselineskip}
\textbf{（二）小位移展开与周期}

在 \(\delta=0\) 附近对式（8）展开到一阶：

\[
\alpha(\delta)
=\alpha(0)+\alpha'(0)\delta+O(\delta^2). \tag{9}
\]

逐项求导可得

\[
\alpha'(0)=\frac{9\sqrt3}{2}. \tag{10}
\]

因此

\[
\alpha(\delta)=\frac{9\sqrt3}{2}\delta+O(\delta^2). \tag{11}
\]

这里的 \(F\) 向上为正，而位移 \(x=\delta a\)
取向下为正，所以动力学方程必须写成

\[
m\ddot x=-F. \tag{12}
\]

保留一阶项：

\[
\begin{aligned}
m\ddot x
&=-\frac{kQq}{a^2}\frac{9\sqrt3}{2}\frac xa\\
&=-\frac{9\sqrt3}{2}\frac{kQq}{a^3}x.
\end{aligned} \tag{13}
\]

与 \(\ddot x+\Omega^2x=0\) 比较：

\[
\Omega^2=\frac{9\sqrt3}{2}\frac{kQq}{ma^3}. \tag{14}
\]

于是

\[
\begin{aligned}
T
&=\frac{2\pi}{\Omega}\\
&=2\pi\sqrt{\frac{2}{9\sqrt3}}
\sqrt{\frac{ma^3}{kQq}}.
\end{aligned} \tag{15}
\]

故

\[
\boxed{\beta=2\pi\sqrt{\frac{2}{9\sqrt3}}\simeq2.250574}. \tag{16}
\]

\needspace{3\baselineskip}
\textbf{（三）\(D\) 点的受力}

\(D\) 到 \(B\)、\(C\) 的连线均沿水平方向，所以 \(B\)、\(C\) 对 \(q\)
的竖直力分量为零，水平分量互相抵消。只有 \(A\) 贡献竖直力。

等边三角形高为 \(\sqrt3a/2\)，故

\[
AD=\frac{\sqrt3}{2}a. \tag{17}
\]

同号排斥方向向下，所以

\[
F_D=-\frac{kQq}{(\sqrt3a/2)^2}
=-\frac{4}{3}\frac{kQq}{a^2}. \tag{18}
\]

因此

\[
\boxed{\alpha'=-\frac43}. \tag{19}
\]

\needspace{3\baselineskip}
\textbf{（四）三角形的有限平衡点}

中心O是一个平衡点。以外接圆半径\(R_c=a/\sqrt3\)为单位，沿中心指向一条边中点的射线取\(r\in(0,1/2)\)。三电荷的径向场与
\[G_3(r)=\frac{2(r-1/2)}{(r^2-r+1)^{3/2}}+\frac1{(1+r)^2}\]
成正比。该函数在零点处的条件可写为 \[(1-2r)(1+r)^2=(r^2-r+1)^{3/2}.\]
两侧在该区间皆为正，可平方并约去中心根，得到
\[r^5+5r^4+r^3+r^2-4r+1=0.\]
非零根\(r=0.2847183548\)，相对于边长的位移为\(\delta=r/\sqrt3=0.1643822188\)。多项式的导数在\((0,1/2)\)内只过零一次，原多项式先减后增；在\(r=0\)为正、\(r=1/2\)为负，因此该区间只有一个根。角向场的严格分析把非中心平衡点限制在边中点射线上，三条射线各一个，故
\[\boxed{N_3=4}.\]
图4-2中的三个附加点按真实\(r\)绘制；它们沿径向是势能极大、沿横向是极小，属于鞍点。

\needspace{3\baselineskip}
\textbf{（五）中心在PQ上的稳定性}

顶点方向单位矢量满足\(\sum\boldsymbol n_j\boldsymbol n_j^{\mathsf T}=3\boldsymbol I/2\)。中心处的场雅可比矩阵为
\[\frac{kQ}{R_c^3}\sum(\boldsymbol I-3\boldsymbol n_j\boldsymbol n_j^{\mathsf T})
=-\frac{9\sqrt3}{2}\frac{kQ}{a^3}\boldsymbol I.\]
因而水平小位移\(\xi\)受到
\[F_x=-\frac{9\sqrt3}{2}\frac{kQq}{a^3}\xi+O(\xi^2),\]
中心在该约束方向稳定。这里始终限制在二维平面内；垂直于平面的微扰反而向外，不能把结论称为三维静电稳定。

\needspace{3\baselineskip}
\textbf{（六）正方形}

同样以外接圆半径为1，沿边中点射线可取顶点角度为\(\pi/4,3\pi/4,5\pi/4,7\pi/4\)。径向场为
\[G_4(r)=\sum_{j=0}^3\frac{r-\cos[(2j+1)\pi/4]}
{[1+r^2-2r\cos((2j+1)\pi/4)]^{3/2}}.\]
除中心外，在\(0<r<\cos(\pi/4)\)内只有一个零点\(r=0.5469363751\)。四条边中点射线给出四个鞍点，故
\[\boxed{N_4=5}.\]

\needspace{3\baselineskip}
\textbf{（七）正2020边形：计数的证据范围}

一般正n边形的径向函数为
\[G_n(r)=\sum_{j=0}^{n-1}\frac{r-\cos[(2j+1)\pi/n]}
{[1+r^2-2r\cos((2j+1)\pi/n)]^{3/2}}.\]
中心附近\(G_n'(0)=-n/2<0\)；圆周上\(G_n(1)>0\)。连续性保证每条边中点射线至少有一个非中心零点，结合中心，得到至少\(n+1\)个平衡点。

对\(n=2020\)数值计算，找到边中点射线上的一个根 \[r=0.999994755115226,\]
旋转复制给出2020个非中心点。常用的预期答案是2021；本次扫描与该数目一致。但``至少一个''不能推出``恰好一个''，密集扫描也不是排除所有附加根的严格证明。

因此这一问应区分：
\[\boxed{N_{2020}\ge2021\quad\text{为已证明的下界}},\]
\[\boxed{N_{2020}=2021\quad\text{为预期答案，本讲义尚未完成其唯一性证明。}}\]
原稿仅凭对称性写出\(n+1\)，逻辑缺少一般n的径向唯一性。讲解时不能把这一缺口隐藏成``显然''。

\textbf{所用的严格结果与限制。}
Bilogliadov的研究给出正多边形Riesz势的角向积分表示，把非中心临界点限制在边中点射线上，并证明每条射线存在非中心临界点；其中\(n=3,4\)的库仑情形可确证总数\(n+1\)。文中对全部n证明``恰为\(n+1\)''的特定指数是\(\beta=1\)，对应\(1/r^2\)势，而本题库仑势对应\(\beta=1/2\)，不能混用。参考：\href{https://math.okstate.edu/people/mbilogli/Bilogliadov_Maxwell.pdf}{Equilibria
of the field generated by point charges}。

\textbf{符号与二阶场。}
三角形中线方向的精确\(\alpha(\delta)\)不是奇函数；三重旋转对称保证一阶雅可比各向同性，不保证沿任一射线的全部非线性项均为奇次。求非中心根要使用完整力函数，不能用线性近似延伸到边缘。

\textbf{数值求根的可复现范围。}
在\((0,1)\)划分区间，近圆周处采用对数加密，计算\(G_n\)并定位变号；再二分至给定精度，验证径向残差、横向分量与旋转复制的坐标。对\(n=2020\)的根，\(1-r\simeq5.245\times10^{-6}\)，均匀粗网格会漏掉。该步骤核对存在性和位置，不宣称完成唯一性证明。

\needspace{3\baselineskip}
\textbf{中心频率和三维稳定性分别检查}

在距中心较近处，把第j个正电荷位置写为\(\boldsymbol a_j\)，长度均为\(a/\sqrt3\)。单项库仑势的二阶展开为
\[\frac1{|\boldsymbol a_j-\boldsymbol r|}
=\frac1{a_j}
+\frac{\boldsymbol a_j\cdot\boldsymbol r}{a_j^3}
+\frac{3(\boldsymbol a_j\cdot\boldsymbol r)^2-a_j^2r^2}{2a_j^5}
+O(r^3).\] 三条方向相加后一次项为零，在三角形平面内
\[\sum_j(\boldsymbol a_j\cdot\boldsymbol r)^2
=\frac32\frac{a^2}{3}r^2.\] 因此平面内的二次势能是
\[\Delta U=\frac{9\sqrt3\,kqQ}{4a^3}r^2,\] 给出
\[\omega^2=\frac{9\sqrt3\,kqQ}{2ma^3},\qquad
T=2\pi\sqrt{\frac{2ma^3}{9\sqrt3\,kqQ}}.\]
把\(T\)写成原题\(T=\beta\sqrt{ma^3/(kqQ)}\)后，才得正文的\(\beta\)。若将\(2\pi\)预先放进定义，会改变要求报告的系数。

若允许沿垂直平面方向z运动，势能二次项却是
\[\Delta U_z=-\frac{9\sqrt3\,kqQ}{2a^3}z^2.\]
所以正试探电荷在中心的平面小振动稳定，而垂直方向不稳定。中心处三维拉普拉斯方程使三方向曲率之和为零，这也是独立符号检验。这里的二维稳定依赖题目限制粒子在平面内运动。

\needspace{3\baselineskip}
\textbf{正多边形中心的小振动可直接作一般化}

设 \(N\ge3\) 个等电荷位于外接圆半径 \(R_c\) 上，单位半径向量为
\(\boldsymbol n_j\)。平面内小位移为
\(\boldsymbol r\)，单个电荷的势能展开为
\[\frac{kqQ}{|\boldsymbol R_j-\boldsymbol r|}
=\frac{kqQ}{R_c}
+\frac{kqQ}{R_c^2}\boldsymbol n_j\cdot\boldsymbol r
+\frac{kqQ}{2R_c^3}
[3(\boldsymbol n_j\cdot\boldsymbol r)^2-r^2]+O(r^3).\] 正多边形满足
\(\sum\boldsymbol n_j=0\)、\(\sum\boldsymbol n_j\boldsymbol n_j^{\mathsf T}=N{\bf1}/2\)，故总势能二阶项为
\[U-U_0=\frac{NkqQ}{4R_c^3}r^2+O(r^3).\] 两个平面模式的角频率相同：
\[\omega_{\parallel,N}^2=\frac{NkqQ}{2mR_c^3}.\] 三角形取
\(R_c=a/\sqrt3\)，回到正文 \(\omega^2=9\sqrt3kqQ/(2ma^3)\)；正方形取
\(R_c=a/\sqrt2\)，则
\(\omega^2=4\sqrt2kqQ/(ma^3)\)。这只需要二阶对称张量，完全不必寻找中心之外的平衡点。
法向位移 \(z\) 的二阶势能为
\(-NkqQz^2/(2R_c^3)\)，仍为不稳定模式。一般中心频率的解析公式与一般平衡点总数的上界是两项独立问题，不能用前者的简洁结果替代后者的全域证明。

\subsection{六、关键洞察}\label{ux516dux5173ux952eux6d1eux5bdf-3}

小位移的稳定性由场的线性系数决定；平衡点的全局计数还需要排除其他零点。对称性生成的零点给出下界，只有附加唯一性分析才能把下界写成等号。

\subsection{七、易错点}\label{ux4e03ux6613ux9519ux70b9-3}

周期中的\(\beta\)须按原题没有额外\(2\pi\)的定义报告。二维中心稳定不代表三维稳定。计数时不要把扫描未发现其他根当作严格证明，也不能把不同势函数指数的定理互相套用。

\subsection{八、课堂讲解建议}\label{ux516bux8bfeux5802ux8bb2ux89e3ux5efaux8bae-3}

先从精确三角形场求周期，再用中心的雅可比解释水平稳定。展示三角形和正方形的径向零点；第七问明确列出下界、数值候选数目与待补的唯一性证明。

\subsection{九、变式训练}\label{ux4e5dux53d8ux5f0fux8badux7ec3-3}

若试探电荷与顶点电荷异号，小位移方程的力整体反向，中心由稳定变为不稳定。还可把正三角形改为一般等腰三角形，要求分别求中心附近纵、横两个方向的线性力系数，观察各向同性何时消失。

\needspace{3\baselineskip}
\textbf{2020边形计数的证据分层}

本题可严格使用旋转对称性：每条相邻电荷角平分线上，一旦证明存在一个非中心零点，其旋转像给出2020个不同零点；再加中心，得到至少2021个。数值根只帮助定位半径，不负责排除额外的径向根。

以单位外接半径写每条角平分线上的径向函数 \[g_n(r)=\sum_{j=0}^{n-1}
\frac{r-\cos[(2j+1)\pi/n]}{[1+r^2-2r\cos((2j+1)\pi/n)]^{3/2}}.\]
中心附近\(g_n'(0)=-n/2<0\)，而\(r=1\)每项非负且总和为正，连续性保证\((0,1)\)至少有一个根。

要把``至少''升级为``恰好''，还须证明角向零点只能位于这些对称射线，并证明每条射线恰有一个正根。当前材料没有为\(n=2020\)完成后一项排他性证明。因此教师版保留``2021为候选总数、严格下界为2021''的标记，不要求学生用图形或扫描替代证明。引用势函数指数不相同的定理，不能补上这个缺口。

\needspace{3\baselineskip}
\textbf{中心在三维空间中的稳定性}

三电荷所在平面内，设试探电荷偏移 \((x,y)\)。在三角形中心附近，二阶势能为
\[U-U_0=\frac{9\sqrt3\,kqQ}{4a^3}(x^2+y^2)
-\frac{9\sqrt3\,kqQ}{2a^3}z^2+O(|\boldsymbol r|^3).\]
因此平面内两个模式具有相同角频率
\[\omega_\parallel^2=\frac{9\sqrt3\,kqQ}{2ma^3},\] 而法向模式满足
\[\ddot z-\frac{9\sqrt3\,kqQ}{ma^3}z=0.\]
正负二次系数之和为零，是无电荷区域 \(\nabla^2U=0\)
的局部检查。原题的稳定小振动必须依赖``限制在平面内''；撤去限制后中心是鞍点。

法向初值 \(z_0,\dot z_0\) 的解为
\[z(t)=z_0\cosh(\sqrt2\omega_\parallel t)
+\frac{\dot z_0}{\sqrt2\omega_\parallel}
\sinh(\sqrt2\omega_\parallel t).\] 增长率是平面振动角频率的 \(\sqrt2\)
倍。学生若从平面中心恢复力推断三维稳定，可以用这个解直接反驳。

对于2020边形，中心外各中线存在数值找到的平衡点，且可由连续性构造中心加各中线的至少2021个。要断言``恰有2021个''，还须排除离开中线的额外点并证明每条中线唯一。低阶三角形、正方形的证明及数值扫描并不代替这两项一般证明。

\newpage

\section{第 5 题}\label{ux7b2c-5-ux9898}

用紧绕硬币的线圈与充电电容器构成放电回路，可在极短时间内使金属硬币受到很大的径向电磁力。建立如下理想模型：线圈与硬币的每一环形细分所产生的磁场在其内部均匀、外部为零；忽略硬币本身的自感。电容为
\(C\)，初始电荷量为 \(Q_0\)；线圈半径为 \(R\)、匝数为 \(N\)、轴向厚度为
\(d\)；硬币电阻率为 \(\rho\)、半径也为 \(R\)，相对磁导率
\(\mu_r\simeq1\)；放电回路的等效电阻为
\(R_0\)。若回路方程出现高阶特征方程，可以用其根表示结果。求时刻 \(t\)
硬币内部单位体积所受线圈磁场力随半径 \(r\) 的分布。

\begin{figure}[H]
\centering
\includegraphics[width=0.67\linewidth,height=\textheight,keepaspectratio,alt={图5-1　电容、线圈和金属硬币组成的脉冲放电装置}]{../figures/question-05.pdf}
\caption*{图5-1　电容、线圈和金属硬币组成的脉冲放电装置}
\end{figure}

\subsection{一、题型定位}\label{ux4e00ux9898ux578bux5b9aux4f4d-4}

线圈电流使硬币产生涡流，涡流又向线圈反馈磁通。体力的径向分布和电流的时间函数分别计算；原题的零自感理想化还要检查是否保持无源性。

\subsection{二、知识点}\label{ux4e8cux77e5ux8bc6ux70b9-4}

法拉第定律；局域欧姆定律；环形电流与磁通链；受涡流反馈的回路方程；三阶特征多项式；体积力密度。

\subsection{三、物理图像}\label{ux4e09ux7269ux7406ux56feux50cf-4}

取轴向磁场正向为\(+z\)，环向电流正向为逆时针。线圈场\(B_0\)增强时，硬币电流沿负环向，故所受线圈磁力指向圆心。

\begin{figure}[H]
\centering
\includegraphics[width=0.798\linewidth,height=\textheight,keepaspectratio,alt={图5-2　环形微元反馈磁通与径向体力}]{../figures/solution-05.pdf}
\caption*{图5-2　环形微元反馈磁通与径向体力}
\end{figure}

半径\(r\)、宽度\(\mathrm dr\)的微元电流为\(\mathrm di=J_\varphi d\,\mathrm dr\)。它产生的磁场按题设只存在于该环以内；线圈圈面积中只有\(\pi r^2\)计入该微元的反馈磁通。

\subsection{四、解题入口}\label{ux56dbux89e3ux9898ux5165ux53e3-4}

记线圈电流为\(I\)，电容剩余电荷为\(Q\)，\(\dot Q=-I\)。先求涡流\(J_\varphi\)，再计算所有环形微元在N匝线圈中的总磁通链。回路电动势由总磁通链的时间导数给出。

\subsection{五、解析}\label{ux4e94ux89e3ux6790-4}

\needspace{3\baselineskip}
\textbf{（一）给定线圈电流时的空间分布}

均匀线圈场及自感为
\[B_0=\frac{\mu_0N}{d}I,\qquad L=\frac{\mu_0N^2\pi R^2}{d}.\]
忽略硬币自感的约定，使每个环形微元的感应电场只取线圈场的变化。半径r圆周的法拉第定律给出
\[2\pi rE_\varphi=-\pi r^2\dot B_0,\]
\[\boxed{J_\varphi=-\frac r{2\rho}\dot B_0}.\]
题目问线圈磁场对硬币的体力，故使用\(\boldsymbol J\times\boldsymbol B_0\)：
\[\boxed{\boldsymbol f(r,t)=-\frac r{2\rho}B_0\dot B_0\,\boldsymbol e_r
=-\frac{\mu_0^2N^2r}{2\rho d^2}I\dot I\,\boldsymbol e_r}.\]
这个空间关系在固定时间对r为线性，轴心处为零。

\needspace{3\baselineskip}
\textbf{（二）硬币电流反馈到线圈}

环形微元电流为\(\mathrm di=J_\varphi d\,\mathrm dr\)。按题设，该微元在自己的内部产生均匀场
\[\mathrm dB_{\rm coin}=\frac{\mu_0\,\mathrm di}{d}=\mu_0J_\varphi\,\mathrm dr.\]
线圈每一匝所链磁通为\(\pi r^2\mathrm dB_{\rm coin}\)，不是\(\pi R^2\mathrm dB_{\rm coin}\)。所以N匝总反馈磁通链
\[\begin{aligned}
\Phi_{\rm coin}&=N\int_0^R\pi r^2\mu_0J_\varphi\,\mathrm dr\\
&=-\frac{\mu_0^2N^2\pi}{2\rho d}\dot I\int_0^Rr^3\,\mathrm dr\\
&=-K\dot I,
\end{aligned}\] 其中 \[\boxed{K=\frac{\mu_0^2N^2\pi R^4}{8\rho d}}.\]
图5-2左半图把\(\pi r^2\)与\(\pi R^2\)分开，正是反馈积分中四次幂\(R^4\)的来源。

\needspace{3\baselineskip}
\textbf{（三）三阶回路方程}

总磁通链为\(LI-K\dot I\)，故回路式是
\[\boxed{L\dot I-K\ddot I+R_0I=Q/C}.\] 对时间求导，并代入\(\dot Q=-I\)：
\[\boxed{-K I^{(3)}+L\ddot I+R_0\dot I+I/C=0}.\]
这解释了原题允许高阶特征方程。忽略\(K\)后才回到孤立线圈的二阶RLC，不能在保留硬币反馈的情况下直接用那个二阶电流。

\needspace{3\baselineskip}
\textbf{（四）特征根形式解及初态}

若电流和涡流初始均为零，则\(I(0)=0\)；由\(J_\varphi\propto\dot I\)又有\(\dot I(0)=0\)。回路式在\(t=0\)给出
\[\ddot I(0)=-\frac{Q_0}{CK}.\] 定义 \[P(s)=-Ks^3+Ls^2+R_0s+1/C.\]
若三个根\(s_j\)互异，拉普拉斯变换为 \[\widehat I(s)=\frac{Q_0}{CP(s)},\]
故形式解
\[\boxed{I(t)=\frac{Q_0}{C}\sum_{j=1}^3\frac{e^{s_jt}}{P'(s_j)}}.\]
有重根时取此表达式的连续极限，出现\(t e^{st}\)或\(t^2e^{st}\)。将该\(I\)及其导数代入第一部分的体力式，即得到原模型所要求的特征根答案。

若题设没有要求初始涡流为零，则应保留\(\dot I(0)\)作为第三个初始数据，不能凭\(I(0),Q(0)\)独自确定三阶方程的解。

\needspace{3\baselineskip}
\textbf{（五）形式模型的物理矛盾}

\(P(0)=1/C>0\)，而\(P(s)\to-\infty\)当\(s\to+\infty\)；故至少有一个正实根。上述无初始电流的解一般包含指数增长，且起初\(\ddot I<0\)，出现与普通放电相反的初始建立方向。这不是硬币获得额外能源，而是``互感非零、硬币自感为零''这一理想化失去无源性。

能量矩阵对一组耦合电感要求\(M^2\le L L_{\rm coin}\)。若\(L_{\rm coin}=0\)却保留\(M\ne0\)，磁能二次型不再非负。原题模型不能作为实际缩币装置的稳定时域预测；应保留硬币自感及磁扩散，或明确采用涡流弱反馈近似。

\textbf{附加弱反馈近似。}
若时间尺度为\(\tau\)且\(K/(L\tau)\ll1\)，只在受迫慢变分支上可以忽略\(K\ddot I\)，用
\[L\ddot I+R_0\dot I+I/C=0.\] 其两根\(\lambda_{1,2}\)给出
\[I_{\rm weak}=\frac{Q_0}{LC(\lambda_1-\lambda_2)}
(e^{\lambda_1t}-e^{\lambda_2t}).\]
此式应标为附加近似，不能与原三阶模型的零涡流初态同时当作精确解。把高阶小项删掉是奇异摄动，增长模态和初始层必须另作处理。

\textbf{力与合力。}
当\(B_0\dot B_0>0\)，力密度向内；反之向外。圆盘的矢量总力因轴对称为零，但径向应力不为零。面积平均的力密度大小为边缘值的\(2/3\)，因为权重是\(2\pi r\,\mathrm dr\)。

\textbf{反馈系数量纲。}
\(K\)的量纲为电感乘时间；\(K\ddot I\)与\(L\dot I\)均为电压。把微元磁通面积错写成整圈\(\pi R^2\)会造成系数和半径幂次同时错误。

\needspace{3\baselineskip}
\textbf{初值和特征根系数如何相互核对}

互异根表达式中的权重\(1/P'(s_j)\)满足部分分式恒等式
\[\frac1{P(s)}=\sum_j\frac1{P'(s_j)(s-s_j)}.\]
两边在大s处展开，三次多项式的首项为\(-Ks^3\)，故
\[\sum_j\frac1{P'(s_j)}=0,\quad
\sum_j\frac{s_j}{P'(s_j)}=0,\quad
\sum_j\frac{s_j^2}{P'(s_j)}=-\frac1K.\]
乘上\(Q_0/C\)后，恰好恢复\(I(0)=0\)、\(\dot I(0)=0\)、\(\ddot I(0)=-Q_0/(CK)\)。因此三个模态都要保留；删掉正增长根一般会破坏给定初值，而不能仍称为原模型的精确解。

再从体积积分检查涡流的焦耳功率：
\[P_{\rm J}=\int_0^R\rho J_\varphi^2\,2\pi rd\,\mathrm dr
=\frac{\pi dR^4}{8\rho}\dot B_0^2.\] 代入\(\dot B_0=\mu_0N\dot I/d\)，有
\[P_{\rm J}=K\dot I^2\ge0.\] 将原回路式乘I，得到
\[\frac{\mathrm d}{\mathrm dt}
\left[\frac{Q^2}{2C}+\frac{LI^2}{2}-KI\dot I\right]
=-R_0I^2-K\dot I^2.\]
右边是非正耗散项，但括号中的``储能''含不定号交叉项\(-KI\dot I\)，无法保证有下界。这个能量式具体展示了零自感模型何处破坏无源性，不能用焦耳功率为正来证明整个模型稳定。

\needspace{3\baselineskip}
\textbf{一个闭合涡流环怎样恢复物理能量账本}

将硬币的连续环流暂用一条等效环代替，设自感为 \(L_c\)、电阻为
\(R_c\)、与线圈互感为 \(M_c\)，环电流为
\(i_c\)。两回路的互感项必须成对保留： \[L\dot I+M_c\dot i_c+R_0I=Q/C,\]
\[M_c\dot I+L_c\dot i_c+R_ci_c=0,\qquad \dot Q=-I.\] 第一式乘
\(I\)，第二式乘 \(i_c\)，再加上电容能量导数，得到
\[\frac{\mathrm d}{\mathrm dt}\left[
\frac{Q^2}{2C}+\frac12LI^2+M_cIi_c+\frac12L_ci_c^2
\right]=-R_0I^2-R_ci_c^2.\] 只要
\(L>0,L_c>0,M_c^2<LL_c\)，磁能二次型为正，电阻又非负，系统不可能凭自身产生增长模态。

若删去 \(L_c\dot i_c\)，第二式变成
\(i_c=-M_c\dot I/R_c\)。代回第一式，确会产生
\[-\frac{M_c^2}{R_c}\ddot I,\] 与原题形式高阶方程中的 \(-K\ddot I\)
同号。但高频支路上 \(\dot i_c\)
很大，先删环自感的近似恰在这里失效。保留原题的形式答案时，应同时说明它不能作为全过程无源瞬态模型；这不是改动一个符号即可解决的问题。

低频或缓慢变化时，条件 \(|s|L_c\ll R_c\)
可以使近似有效。若要求开启瞬间的完整暂态，则应使用双回路方程和两个初始电流，而不是在降阶近似里追加任意初值。

\subsection{六、关键洞察}\label{ux516dux5173ux952eux6d1eux5bdf-4}

空间上的\(r\)来自磁通面积与环周长之比；时间上的三阶方程来自涡流反馈磁通含有\(\dot I\)。这两个步骤分别决定空间分布与电流时间函数。

\subsection{七、易错点}\label{ux4e03ux6613ux9519ux70b9-4}

不要漏掉硬币向线圈反馈的磁通。微元内部磁场只链面积\(\pi r^2\)。初始涡流为零不能与二阶孤立RLC的\(\dot I(0)=Q_0/(LC)\)混用。正增长根须作为模型缺陷说明。

\subsection{八、课堂讲解建议}\label{ux516bux8bfeux5802ux8bb2ux89e3ux5efaux8bae-4}

先由已知\(B_0(t)\)求体力，再计算反馈磁通并列回路。让学生检查特征多项式在0与正无穷的符号，判断这个简化模型是否仍然无源。

\subsection{九、变式训练}\label{ux4e5dux53d8ux5f0fux8badux7ec3-4}

若线圈场仅存在于\(r<a<R\)，则\(r>a\)处的环向电场是\(E_\varphi=-a^2\dot B_0/(2r)\)。那里虽然有感生电流，但线圈局部磁场为零，所以题目所问的线圈场体力为零；不能仅把内区公式中的\(r\)替换成\(a^2/r\)便宣称外区有同样的磁力。

\needspace{3\baselineskip}
\textbf{保留硬币自感后的无源双回路}

用一个等效硬币回路演示模型修复。记其自感\(L_c\)、电阻\(R_c\)、互感\(M\)、电流\(i_c\)，线圈电流为I。取相容电流方向，方程为
\[L\dot I+M\dot i_c+R_0I=Q/C,\]
\[M\dot I+L_c\dot i_c+R_ci_c=0,\qquad\dot Q=-I.\] 磁能和电容能总量为
\[E=\frac{LI^2}{2}+Mi_cI+\frac{L_ci_c^2}{2}+\frac{Q^2}{2C}.\]
只要\(L>0,L_c>0,M^2<LL_c\)，磁能二次型正定。两回路分别乘相应电流相加，得到
\[\dot E=-R_0I^2-R_ci_c^2\le0.\]
由此排除无外源的持续指数增长。真实硬币需要连续环形回路或磁扩散模型，但这一双回路例已足以说明为什么不能一面保留互感、一面令硬币自感严格为零。

把第二回路形式上近似为\(i_c=-M\dot I/R_c\)，会在第一回路产生\(-M^2\ddot I/R_c\)。这个消元只适用于变化比\(L_c/R_c\)慢的分支，并非对所有特征根和初态都一致有效。

\needspace{3\baselineskip}
\textbf{对称双回路的正定能量例子}

令 \(L_c=L\)、\(M_c=L/2\)，两回路电阻均为
\(R\)，暂不接电容。自然衰减方程为 \[L\dot I+\frac L2\dot i_c+RI=0,\]
\[\frac L2\dot I+L\dot i_c+Ri_c=0.\] 取
\(I_+=I+i_c,I_-=I-i_c\)，相加与相减后得到
\[\frac{3L}{2}\dot I_++RI_+=0,\qquad
\frac L2\dot I_-+RI_-=0.\] 所以两种模式分别按
\[I_+(t)=I_+(0)e^{-2Rt/(3L)},\qquad
I_-(t)=I_-(0)e^{-2Rt/L}\] 衰减。磁能可以写成
\[U_B=\frac{3L}{8}I_+^2+\frac L8I_-^2>0.\]
两项都为正，与两个衰减率同号。若在这组方程中把 \(L_c\)
单独删去而仍保留互感，磁能矩阵行列式变成
\(-M_c^2<0\)，就会引入非物理分支。
恢复电容后系统可有阻尼振荡，但总储能的时间导数仍是两只电阻的负耗散功率。这一结论无需逐个求复杂特征根，更适合判断原题形式高阶方程能否描述完整瞬态。

\needspace{3\baselineskip}
\textbf{形式高阶方程的长时间分支如何筛选}

原题三阶方程的特征多项式为 \[P(s)=-Ks^3+Ls^2+R_0s+1/C.\]
对正实轴，\(P(0)>0\) 而
\(P(s)\to-\infty\)，因此至少有一个正实根。即使初始电流为零，电容初始电荷会通过初始导数条件激发这个分支；不能凭
\(I(0)=0\) 宣称增长项系数为零。
若仅考虑慢变近似，须从保留硬币自感的完整系统中选择随 \(L_c\to0\)
连续趋于有限慢模态的支路，而不是在已经病态的三阶方程里事后任意删去增长指数。完整系统的磁能正定条件为
\(M_c^2<LL_c\)；把 \(L_c\) 趋零而 \(M_c\)
固定也违反这一条件，因此参数极限必须一并交代。
用于评分时可分别记录两项：按原题删自感模型推出正确的形式方程；识别该方程作为完整瞬态模型的局限。不能用普通二阶RLC方程替换第一项，又不说明删去了硬币反馈。

\newpage

\section{第 6 题}\label{ux7b2c-6-ux9898}

在地面惯性系 \(S\) 中，猎人静止在原点 \(O\)，小猪起初静止在 \(A\)
点，\(OA=x_1\)。\(t=0\) 时猎人沿 \(+x\) 方向发射激光，小猪同时以恒定速度
\(v\) 沿 \(+x\) 方向逃跑。激光在事件 \(E_2(t_2,x_2)\)
追上小猪，随后反射光返回并在事件 \(E_3(t_3,0)\) 到达猎人。另一惯性系
\(S'\) 以速度 \(u\) 沿 \(+x\) 方向相对 \(S\) 运动，且两坐标系原点在
\(t=t'=0\) 重合。求：

（1）地面系中的 \(t_3\)；

（2）事件 \(E_3\) 在 \(S'\) 中的时刻 \(t_3'\)；

（3）事件 \(E_3\) 在 \(S'\) 中的位置 \(x_3'\)；

（4）小猪起跑事件 \(E_1(0,x_1)\) 在 \(S'\) 中的时刻 \(t_1'\)；

（5）小猪被击中事件在 \(S'\) 中的位置 \(x_2'\)。

\begin{figure}[H]
\centering
\includegraphics[width=0.78\linewidth,height=\textheight,keepaspectratio,alt={图6-1　地面系中猎人、激光与匀速逃跑小猪的初始关系}]{../figures/question-06.pdf}
\caption*{图6-1　地面系中猎人、激光与匀速逃跑小猪的初始关系}
\end{figure}

\subsection{一、题型定位}\label{ux4e00ux9898ux578bux5b9aux4f4d-5}

把起跑、击中、返回分别看成事件。先在地面系求它们的坐标，再逐个作洛伦兹变换；回程光速用事件差复核。

\subsection{二、知识点}\label{ux4e8cux77e5ux8bc6ux70b9-5}

\begin{enumerate}
\def\labelenumi{\arabic{enumi}.}
\tightlist
\item
  光速不变与一维追及方程。
\item
  光的往返传播时间。
\item
  洛伦兹变换 \(t'=\gamma_u(t-ux/c^2)\)、\(x'=\gamma_u(x-ut)\)。
\item
  同时性的相对性。
\item
  同地事件的固有时与非同地事件的坐标时区别。
\end{enumerate}

\subsection{三、物理图像}\label{ux4e09ux7269ux7406ux56feux50cf-5}

在 \(x\)-\(ct\) 图中，出射光和反射光的世界线斜率绝对值均为
\(1\)；小猪世界线从 \((x_1,0)\) 出发，斜率由 \(v\) 决定；猎人世界线是
\(x=0\)。

\begin{figure}[H]
\centering
\includegraphics[width=0.795\linewidth,height=\textheight,keepaspectratio,alt={图6-2　起跑、击中及返回事件的洛伦兹世界线}]{../figures/solution-06.pdf}
\caption*{图6-2　起跑、击中及返回事件的洛伦兹世界线}
\end{figure}

图中三个事件分别为

\[
E_1=(0,x_1),\qquad E_2=(t_2,x_2),\qquad E_3=(t_3,0). \tag{1}
\]

\subsection{四、解题入口}\label{ux56dbux89e3ux9898ux5165ux53e3-5}

定义

\[
\gamma_u=\frac1{\sqrt{1-u^2/c^2}}. \tag{2}
\]

标准洛伦兹变换为

\[
\boxed{
t'=\gamma_u\left(t-\frac{ux}{c^2}\right),\qquad
x'=\gamma_u(x-ut)}. \tag{3}
\]

由于两系原点在发射事件重合，不需要添加额外的时间或空间平移常量。

\subsection{五、解析}\label{ux4e94ux89e3ux6790-5}

\needspace{3\baselineskip}
\textbf{（一）激光追上小猪的事件 \(E_2\)}

在地面系中，激光位置为

\[
x_{\rm light}=ct. \tag{4}
\]

小猪位置为

\[
x_{\rm pig}=x_1+vt. \tag{5}
\]

追上时二者位置相同：

\[
ct_2=x_1+vt_2. \tag{6}
\]

把含 \(t_2\) 的项移到左边：

\[
(c-v)t_2=x_1. \tag{7}
\]

所以

\[
\boxed{t_2=\frac{x_1}{c-v}}. \tag{8}
\]

击中位置为

\[
\boxed{x_2=ct_2=\frac{cx_1}{c-v}}. \tag{9}
\]

\needspace{3\baselineskip}
\textbf{（二）反射光返回猎人的时刻 \(t_3\)}

从 \(E_2\) 到 \(E_3\)，反射光向左走过距离 \(x_2\)，所需时间为

\[
t_3-t_2=\frac{x_2}{c}. \tag{10}
\]

由 \(x_2=ct_2\)，

\[
t_3-t_2=t_2. \tag{11}
\]

于是

\[
t_3=2t_2. \tag{12}
\]

代入式（8）：

\[
\boxed{t_3=\frac{2x_1}{c-v}}. \tag{13}
\]

\needspace{3\baselineskip}
\textbf{（三）\(E_3\) 在 \(S'\) 中的时空坐标}

\(E_3\) 的地面系坐标是 \((t_3,0)\)。代入式（3）：

\[
\begin{aligned}
t_3'
&=\gamma_u\left(t_3-\frac{u\cdot0}{c^2}\right)\\
&=\gamma_ut_3\\
&=\frac{2\gamma_ux_1}{c-v}.
\end{aligned} \tag{14}
\]

因此

\[
\boxed{t_3'=\frac{2\gamma_ux_1}{c-v}}. \tag{15}
\]

位置为

\[
\begin{aligned}
x_3'
&=\gamma_u(0-ut_3)\\
&=-\gamma_uut_3\\
&=-\frac{2\gamma_uux_1}{c-v}.
\end{aligned} \tag{16}
\]

故

\[
\boxed{x_3'=-\frac{2\gamma_uux_1}{c-v}}. \tag{17}
\]

这说明猎人在 \(S'\) 中以速度 \(-u\) 运动，时刻 \(t_3'\) 时恰位于
\(x_3'=-ut_3'\)，与式（15）、（17）完全一致。

\needspace{3\baselineskip}
\textbf{（四）小猪起跑事件的 \(t_1'\)}

\(E_1\) 在地面系中为 \((0,x_1)\)。代入时间变换：

\[
\begin{aligned}
t_1'
&=\gamma_u\left(0-\frac{ux_1}{c^2}\right)\\
&=-\frac{\gamma_uux_1}{c^2}.
\end{aligned} \tag{18}
\]

所以

\[
\boxed{t_1'=-\frac{\gamma_uux_1}{c^2}}. \tag{19}
\]

负号并非矛盾，而是同时性的相对性：在 \(S\)
中猎人发射与小猪起跑同时、但异地；在向右运动的 \(S'\)
中，右侧的小猪起跑事件发生得更早。

\needspace{3\baselineskip}
\textbf{（五）击中事件的 \(x_2'\)}

把式（8）、（9）代入位置变换：

\[
\begin{aligned}
x_2'
&=\gamma_u(x_2-ut_2)\\
&=\gamma_u(c-u)t_2\\
&=\gamma_u(c-u)\frac{x_1}{c-v}.
\end{aligned} \tag{20}
\]

因此

\[
\boxed{x_2'=\frac{\gamma_u(c-u)x_1}{c-v}}. \tag{21}
\]

可以用光世界线作强核验。出射光满足 \(x=ct\)，对其任一事件作洛伦兹变换：

\[
\frac{x'}{t'}=
\frac{x-ut}{t-ux/c^2}
=\frac{c-u}{1-u/c}=c.
\]

因此在 \(S'\) 中光速仍为 \(c\)。同时，\(E_2\) 与 \(E_3\)
的间隔为光样间隔，洛伦兹变换后也必须满足 \(|x_3'-x_2'|=c(t_3'-t_2')\)。

式（15）可以看作时间膨胀，是因为发射事件与返回事件都发生在猎人处，在
\(S\) 中同地；而 \(E_1\) 与发射事件异地，绝不能把它们的时间差用单纯的
\(\gamma_u\) 关系处理。

补算击中事件的时间坐标可形成完整事件表：

\[
\begin{aligned}
t_2'
&=\gamma_u\left(t_2-\frac{ux_2}{c^2}\right)\\
&=\gamma_ut_2\left(1-\frac uc\right)\\
&=\frac{\gamma_u(c-u)x_1}{c(c-v)}.
\end{aligned}
\]

由式（21）立即得到 \(x_2'=ct_2'\)，验证出射光在 \(S'\) 中仍满足
\(x'=ct'\)。反射段则有

\[
x_2'-x_3'=c(t_3'-t_2'),
\]

把式（15）、（17）、（21）和上式 \(t_2'\) 逐项代入，两边都化为

\[
\frac{\gamma_u(c+u)x_1}{c-v}.
\]

此式分别从空间坐标差与回程时间差得到。这个代数核验同时检查了返回光方向、\(x_3'\)
的负号以及两个事件的先后顺序。

还可从两个事件的坐标差恢复小猪在 \(S'\) 中的速度。由 \(E_1\)、\(E_2\)，

\[
v'=\frac{x_2'-x_1'}{t_2'-t_1'}.
\]

把两事件分别作洛伦兹变换并提取公共因子，得到标准速度变换

\[
\boxed{v'=\frac{v-u}{1-uv/c^2}}.
\]

若 \(u=v\)，\(S'\) 与小猪同速，因而 \(v'=0\)；若
\(v=0\)，原来静止的小猪在 \(S'\) 中以 \(-u\) 运动；若令 \(v\to c\)，则
\(v'\to c\)。三个极限分别核查了共动系、静止物体和光速不变性。

还可检查发射事件 \(E_0=(0,0)\) 与返回事件 \(E_3=(t_3,0)\) 的时空间隔：

\[
c^2\Delta\tau^2=c^2t_3^2.
\]

在 \(S'\) 中，

\[
\begin{aligned}
c^2t_3'^2-x_3'^2
&=\gamma_u^2t_3^2(c^2-u^2)\\
&=c^2t_3^2.
\end{aligned}
\]

所以猎人腕表记录的固有时就是 \(t_3\)，而 \(S'\) 的坐标时为
\(t_3'=\gamma_ut_3\)。这项不变量检查比单看公式的 \(\gamma_u\)
因子更能说明为什么第（2）问可使用时间膨胀。

各结果在 \(u=0\) 时必须恢复地面系：\(\gamma_u=1\)，故

\[
t_3'=t_3\qquad x_3'=0\qquad t_1'=0\qquad x_2'=x_2.
\]

当 \(u>0\) 时，\(t_1'<0\)、\(x_3'<0\)；两个负号分别反映右侧异地事件在
\(S'\) 中提前，以及猎人在 \(S'\)
中向左运动。若只得到数值大小而没有这些符号，读者无法从结果重建事件图景。相反，把符号和世界线相互核验后，不必额外背诵``哪个参考系看谁向哪边运动''的口头规则。

\textbf{反射段的独立复核。} 相应两事件的坐标差应满足
\[x_2'-x_3'=\gamma(c+u)\frac{x_1}{c-v}
=c(t_3'-t_2').\]
这个条件检验回程光线，而不是要求所有世界线图中的投影斜率相同。绘图采用相同的\(x'\)与\(ct'\)单位，光线才与坐标轴成\(45^\circ\)。两次同时截面也应由同一洛伦兹时间变换确定，不能凭图中水平线替代。

\needspace{3\baselineskip}
\textbf{在新系直接追及，复核事件表}

小猪在\(S'\)中的速度为\(v'=(v-u)/(1-uv/c^2)\)，猎人速度为\(-u\)。但小猪在\(t'=0\)的位置不是起跑事件的\(x_1'=\gamma_ux_1\)。先从其世界线\(x=x_1+vt\)与\(t'=0\)联立：
\[t=\frac{ux_1}{c^2-uv},\qquad
x_{\rm pig}'(0)=\frac{x_1}{\gamma_u(1-uv/c^2)}.\]
于是出射光\(x'=ct'\)与小猪的相遇时间是
\[t_2'=\frac{x_{\rm pig}'(0)}{c-v'}
=\frac{\gamma_u(c-u)x_1}{c(c-v)},\] 与事件变换相同。

从击中事件开始，回程光位置为\(x_2'-c(t'-t_2')\)，猎人位置为\(-ut'\)。二者相等给出
\[t_3'=\frac{x_2'+ct_2'}{c-u}
=\frac{2\gamma_ux_1}{c-v}.\]
这里同时用了新系的光速、猪速和猎人速度；把猎人固定在\(x'=0\)会漏掉返回过程中的追及距离。图6-2中猎人的世界线随变换倾斜，正是这项计算的几何表示。

作为数值例，取\(u=0.2c,v=0.4c,x_1/c=1\)个时间单位，得到\(t_2=5/3\)、\(t_3=10/3\)，而\(t_2'=1.3608276\)、\(t_3'=3.4020691\)。两个新系时间并不再相差恰好一倍，原因是猎人在回程期间继续向左运动。

\subsection{六、关键洞察}\label{ux516dux5173ux952eux6d1eux5bdf-5}

先把事件在 \(S\)
中求完整，再逐个代入同一套洛伦兹变换。相对论题中``事件表''比零散公式更可靠：每一行都必须同时有
\((t,x)\)，缺少一个坐标便无法变换另一个。

\subsection{七、易错点}\label{ux4e03ux6613ux9519ux70b9-5}

\begin{enumerate}
\def\labelenumi{\arabic{enumi}.}
\tightlist
\item
  追及方程误写为 \(ct_2=x_1-vt_2\)。
\item
  反射返回时间漏掉从 \(E_2\) 开始，误写成 \(t_3=x_2/c\)。
\item
  对异地同时事件直接使用 \(t'=\gamma t\)。
\item
  洛伦兹时间变换漏掉 \(ux/c^2\)。
\item
  认为猎人在任何参考系中都位于原点，因而误写 \(x_3'=0\)。
\end{enumerate}

\subsection{八、课堂讲解建议}\label{ux516bux8bfeux5802ux8bb2ux89e3ux5efaux8bae-5}

让学生先画三条世界线，并把 \(E_1\)、\(E_2\)、\(E_3\)
的地面系坐标写在图旁。之后每问只做``从事件表取一行---代入变换---化简''三步。讲到
\(t_1'<0\) 时，可借图解释同时线倾斜，而不是只说``相对性如此''。

\subsection{九、变式训练}\label{ux4e5dux53d8ux5f0fux8badux7ec3-5}

可要求补求 \(t_2'\)，再在 \(S'\) 中用 \(E_2'\)、\(E_3'\)
直接验证反射光速度为 \(-c\)。也可令小猪向猎人方向运动，即把 \(v\)
改为负值，检查所有公式是否自然退化而无需重新记忆追及公式。

\needspace{3\baselineskip}
\textbf{从激光往返数据反求速度}

若已知地面系初始距离\(x_1\)及猎人腕表往返时间\(t_3\)，由\(t_3=2x_1/(c-v)\)可反求
\[v=c-\frac{2x_1}{t_3}.\]
对向外逃跑的小猪，要求\(t_3>2x_1/c\)；等号对应\(v=0\)。测量误差传播为
\[\delta v=\frac{2x_1}{t_3^2}\delta t_3-\frac2{t_3}\delta x_1.\]
若用的是运动参考系坐标时间\(t_3'\)，必须先除\(\gamma_u\)还原猎人固有时，不能直接代入上式。

猎人雷达定义的击中时刻和距离为\((t_0+t_3)/2\)、\(c(t_3-t_0)/2\)。因\(t_0=0\)且猎人静止，分别为\(t_2\)和\(x_2\)。这组雷达坐标属于猎人的惯性系；对S'中的移动猎人，不能把坐标时一分为二就声称得到S'原点到目标的距离。

\needspace{3\baselineskip}
\textbf{负相对速度的参考系也可直接追及}

取 \(v=0.4c,u=-0.2c,x_1=l\)。新系小猪速度为
\[v'=\frac{0.4c+0.2c}{1+0.08}=\frac59c,\] 猎人速度为
\(-u=+0.2c\)。在新系 \(t'=0\) 截面，小猪位置为
\[x'_1=\frac{l}{\gamma(1-uv/c^2)}
=\frac{\sqrt{0.96}}{1.08}l.\] 出射光追到小猪用时
\[t'_2=\frac{x'_1}{c-v'}=\frac{5\gamma}{3}\frac{1.2l}{c}.\]
这也等于原系事件变换 \(\gamma(t_2-ux_2/c^2)\)，其中
\(t_2=5l/(3c),x_2=5l/3\)。 返回光与向右运动的猎人相向接近，闭合速度为
\(c+0.2c\)，故回程用时为碰面时两者间距除以 \(1.2c\)。将两段时间相加得到
\[t'_3=\gamma t_3=\frac{10\gamma l}{3c}.\]
发射和返回都在猎人世界线上，猎人钟的固有时间仍为原系
\(10l/(3c)\)。新系坐标时间与它的比值为 \(\gamma\)。该例检查 \(u\)
的符号：若沿用 \(c-u\) 而不结合猎人速度方向，会把相向返回错误写成追赶。

\newpage

\section{第 7 题}\label{ux7b2c-7-ux9898}

两个焦距均为\(f\)的理想薄凸透镜共用光心O。图中的\(f_1,f_2\)表示透镜所在的直线，夹角为\(60^\circ\)；\(f_1\)沿y轴，\(f_2\)与正x轴成\(30^\circ\)。\(A=(-f,f/\sqrt3)\)，故\(OA=2f/\sqrt3\)；\(B=(-\sqrt3f,0)\)位于第一透镜的主光轴上。按原题的理想连续成像模型求：

（1）（2）A先经\(f_1\)再经\(f_2\)的终像横、纵坐标\(x_1,y_1\)；

（3）（4）A先经\(f_2\)再经\(f_1\)的终像横、纵坐标\(x_1',y_1'\)；

（5）（6）B先经\(f_1\)再经\(f_2\)的终像横、纵坐标\(x_2,y_2\)；

（7）（8）B先经\(f_2\)再经\(f_1\)的终像横、纵坐标\(x_2',y_2'\)。

\begin{figure}[H]
\centering
\includegraphics[width=0.587\linewidth,height=\textheight,keepaspectratio,alt={图7-1　两条透镜平面、法向与原题物点}]{../figures/question-07.pdf}
\caption*{图7-1　两条透镜平面、法向与原题物点}
\end{figure}

\subsection{一、题型定位}\label{ux4e00ux9898ux578bux5b9aux4f4d-6}

原图的长直线表示透镜平面，物距沿法向取带符号投影。中间像在无穷远时，用齐次坐标继续第二次成像。

\subsection{二、知识点}\label{ux4e8cux77e5ux8bc6ux70b9-6}

透镜平面法向；实物、虚物的带符号物距；过光心光线共线；连续薄透镜成像；齐次坐标中的无穷远像。

\subsection{三、物理图像}\label{ux4e09ux7269ux7406ux56feux50cf-6}

选择光线前进方向的两条法向为
\[\boldsymbol n_1=(1,0),\qquad \boldsymbol n_2=(1/2,-\sqrt3/2).\]
这两向量垂直于图中的\(f_1,f_2\)。A对两法向的投影均为\(-f\)，位于两个透镜的前焦平面。

\begin{figure}[H]
\centering
\includegraphics[width=0.87\linewidth,height=\textheight,keepaspectratio,alt={图7-2　两镜法向、无穷远中间像与终像}]{../figures/solution-07.pdf}
\caption*{图7-2　两镜法向、无穷远中间像与终像}
\end{figure}

\subsection{四、解题入口}\label{ux56dbux89e3ux9898ux5165ux53e3-6}

令\(\boldsymbol r\)是相对光心的位置，\(z=\boldsymbol n\cdot\boldsymbol r\)为带符号轴向坐标，物距\(s=-z\)。薄透镜式及过光心光线给出
\[\boxed{\mathcal L_{\boldsymbol n}(\boldsymbol r)=
\frac{\boldsymbol r}{1+\boldsymbol n\cdot\boldsymbol r/f}}.\]
该式使用有向投影，不能将点积取绝对值，否则虚物的符号会被抹去。

\subsection{五、解析}\label{ux4e94ux89e3ux6790-6}

\needspace{3\baselineskip}
\textbf{（一）单透镜映射的推导}

若\(s\ne f\)，像距\(s'=fs/(s-f)\)，横向放大率\(-s'/s\)。过光心的光线使像点与物点共线，故
\[\boldsymbol r'=-\frac{s'}s\boldsymbol r
=\frac{f}{f-s}\boldsymbol r
=\frac{\boldsymbol r}{1+z/f}.\]
当中间像成为第二镜的虚物，\(s\)取负值，该式仍适用。若\(s=f\)，光束出射平行；不要把无穷远像当作一个普通有限坐标再代入。

\needspace{3\baselineskip}
\textbf{（二）用齐次坐标处理焦平面}

把点记成\((\boldsymbol p,w)\)，有限点的位置为\(\boldsymbol r=\boldsymbol p/w\)。每次透镜作用保持\(\boldsymbol p\)不变，而
\[w\longmapsto w+\boldsymbol n\cdot\boldsymbol p/f.\]
因此\(w=0\)表示无穷远像，两次作用合成后为 \[\boxed{\boldsymbol r_{12}=
\frac{\boldsymbol r}{1+(\boldsymbol n_1+\boldsymbol n_2)\cdot\boldsymbol r/f}}.\]
两次对同一\(w\)作加法，先后次序在这个共光心理想映射中可交换。这里的结论依赖共光心、固定法向及连续薄透镜模型，不能推广到相隔有限距离的真实透镜组。

\needspace{3\baselineskip}
\textbf{（三）A点的八问中前四问}

对\(\boldsymbol r_A=(-f,f/\sqrt3)\)，有
\[\boldsymbol n_1\cdot\boldsymbol r_A=-f,\qquad
\boldsymbol n_2\cdot\boldsymbol r_A=-f.\]
第一次成像后的齐次分母为0，代表平行光；第二次后分母为\(1-1-1=-1\)。所以
\[\boldsymbol r_{A,12}=\boldsymbol r_{A,21}=-\boldsymbol r_A.\]
逐个坐标为 \[\boxed{x_1=f,\quad y_1=-f/\sqrt3,\quad
x_1'=f,\quad y_1'=-f/\sqrt3}.\]
图7-2标明第一次的平行光与第二次的有限终像，不能再画第一次虚像。

\needspace{3\baselineskip}
\textbf{（四）B点的后四问}

\(\boldsymbol r_B=(-\sqrt3f,0)\)，两投影分别是\(-\sqrt3f\)、\(-\sqrt3f/2\)，故合成分母为\(1-3\sqrt3/2\)。于是
\[\boldsymbol r_{B,12}=\boldsymbol r_{B,21}
=\left(\frac{2\sqrt3}{3\sqrt3-2}f,0\right).\] 有理化可写为
\[\frac{2\sqrt3}{3\sqrt3-2}=\frac{18+4\sqrt3}{23}=1.083834923\ldots.\]
因此 \[\boxed{x_2=x_2'=\frac{18+4\sqrt3}{23}f,\qquad y_2=y_2'=0}.\]

\needspace{3\baselineskip}
\textbf{（五）B点逐次复核}

先经第一镜，\(s_1=\sqrt3f\)，第一像 \[\boldsymbol r_B^{(1)}=
\left(\frac{\sqrt3}{\sqrt3-1}f,0\right).\]
这个中间像对第二镜位于其前进一侧，是虚物，应取
\[s_2=-\boldsymbol n_2\cdot\boldsymbol r_B^{(1)}<0.\]
用负物距再次成像，与上面的合成结果一致。原稿取绝对值使虚物变成实物，因而得出错误的\(2\sqrt3f\)。

\textbf{共线与数值检查。}
A的终像满足\(y/x=-1/\sqrt3\)，B的终像仍在x轴上；两者都满足过光心共线条件。A的\(OA'\)与\(OA\)相等、方向相反。B的总放大因子为\(1/(1-3\sqrt3/2)\)，为负，终像在物点另一侧。

\textbf{模型范围。}
大夹角交叉透镜不能视为真实透镜系统中的统一小角近轴传播；原题是指定``仍满足理想薄透镜工作原理''的数学模型。实际光线的穿镜顺序、有限孔径及大角像差需另建光学模型。本题在理想连续成像规则下给出上述坐标，不以它替代实际器材的精确预测。

\textbf{有限距离破坏可交换性。}
两镜之间的传播使齐次位置也发生变化，光线矩阵中的平移项不能省略。共光心映射的相加结构由此消失。讨论次序时应先说明是哪一种模型，而不是把``不对易''当作任何双透镜系统的固定结论。

\needspace{3\baselineskip}
\textbf{反向次序的中间像也要保留符号}

B先经第二镜时，齐次分母为\(1-\sqrt3/2>0\)，故 \[\boldsymbol r_B^{(2)}
=\left(-\frac{\sqrt3}{1-\sqrt3/2}f,0\right).\]
它仍在左侧，且离光心更远。这一中间像与先经第一镜时的右侧中间像不同，但再次经第一镜后
\[\boldsymbol r_B^{(21)}
=\frac{\boldsymbol r_B^{(2)}}
{1+\boldsymbol n_1\cdot\boldsymbol r_B^{(2)}/f}
=\left(\frac{2\sqrt3}{3\sqrt3-2}f,0\right).\]
所以``终像相同''不等于``中间像相同''，也不表示可以把两条路径画成同一条有限光线。

把\(x,y,w\)合成一个齐次列向量，单镜矩阵是 \[H(\boldsymbol n)=
\begin{pmatrix}
1&0&0\\0&1&0\\n_x/f&n_y/f&1
\end{pmatrix}.\] 矩阵相乘给出 \[H(\boldsymbol n_2)H(\boldsymbol n_1)=
\begin{pmatrix}
1&0&0\\0&1&0\\(n_{1x}+n_{2x})/f&(n_{1y}+n_{2y})/f&1
\end{pmatrix}.\]
这个乘积即使经过\(w=0\)也有定义。对A，第一步的齐次像是\((-f,f/\sqrt3,0)\)，第二步成为\((-f,f/\sqrt3,-1)\)；恢复有限坐标后为\((f,-f/\sqrt3)\)。图7-2所标``中间像在无穷远''由此有了可计算的含义。

\needspace{3\baselineskip}
\textbf{物点B的两种次序逐步计算}

把所有坐标除以 \(f\)，令 \(b=\sqrt3\)。原物点为
\(\boldsymbol B=(-b,0)\)。先经过第一镜，其法向点积为 \(-b\)，分母为
\(1-b\)，所以 \[\boldsymbol B_1=\left(\frac{b}{b-1},0\right).\]
这是正横坐标的中间像。第二镜法向的横分量为 \(1/2\)，故
\[1+\boldsymbol n_2\cdot\boldsymbol B_1
=1+\frac{b}{2(b-1)}=\frac{3b-2}{2(b-1)}.\] 于是终像为
\[\boldsymbol B_{12}=\left(\frac{2b}{3b-2},0\right)
=\left(\frac{18+4\sqrt3}{23},0\right).\] 换次序先过第二镜，其分母为
\(1-b/2>0\)，中间像为
\[\boldsymbol B_2=\left(-\frac{2b}{2-b},0\right).\]
这一中间像在负横坐标且距光心较远。再经第一镜，分母为
\[1+\boldsymbol n_1\cdot\boldsymbol B_2
=\frac{2-3b}{2-b}<0.\] 第二次归一化后
\[\boldsymbol B_{21}=\left(\frac{2b}{3b-2},0\right)
=\boldsymbol B_{12}.\]
两个中间像显著不同，终像却相同，故可交换性说的是合成映射。对于A，一种次序中间像在无穷远，要用齐次坐标；不能把``中间像不同''或``中间像无限远''当作终像不一致的证据。

\subsection{六、关键洞察}\label{ux516dux5173ux952eux6d1eux5bdf-6}

图上的透镜平面与主光轴相互垂直。带符号成像保留虚物信息，齐次坐标保留平行光信息，两者都是本题连续成像所必需的。

\subsection{七、易错点}\label{ux4e03ux6613ux9519ux70b9-6}

把\(f_1,f_2\)误当成主光轴、对点积取绝对值、把无穷远像代成有限点，都会改变答案。八个坐标须分别对应四种路径。

\subsection{八、课堂讲解建议}\label{ux516bux8bfeux5802ux8bb2ux89e3ux5efaux8bae-6}

先由原图确认B位于第一镜主光轴，所以第一镜法向沿x轴。再推导带符号映射，利用A的第一次焦平面成像引出齐次表示；最后用B逐次成像复核合成式。

\subsection{九、变式训练}\label{ux4e5dux53d8ux5f0fux8badux7ec3-6}

把两焦距改为\(f_1,f_2\)，在同一共光心理想映射内
\[\boldsymbol r_{12}=\frac{\boldsymbol r}
{1+\boldsymbol n_1\cdot\boldsymbol r/f_1+\boldsymbol n_2\cdot\boldsymbol r/f_2}.\]
终像在无穷远的物点满足一条直线方程。可要求学生在图中画出该``合成前焦平面''，并解释为何它的法向不是任一单镜法向。

\needspace{3\baselineskip}
\textbf{不同焦距与有间距的对照}

共光心模型改为焦距\(f_1,f_2\)后，合成前焦平面为
\[\left(\frac{\boldsymbol n_1}{f_1}
+\frac{\boldsymbol n_2}{f_2}\right)\cdot\boldsymbol r=-1.\]
若\(f_1=f_2=f\)，本题法向和的长度为\(\sqrt3/f\)，这个平面到光心的垂距为\(f/\sqrt3\)，方向与任一单镜都不同。

用普通共轴、有间距d的近轴系统作对照。其光线矩阵为 \[L_2P(d)L_1=
\begin{pmatrix}
1-d/f_1&d\\
-1/f_1-1/f_2+d/(f_1f_2)&1-d/f_2
\end{pmatrix}.\]
交换两透镜后，上左、下右两个元素交换。若\(f_1\ne f_2,d\ne0\)，矩阵一般不同。因此正文的可交换性依赖指定共光心映射，不能把它推广为实际双镜成像的通则。

这个对照使用的是另一套共轴小角模型，不用它直接计算原题的大夹角交叉透镜，而是指出传播距离在哪个数学位置破坏相加结构。

\needspace{3\baselineskip}
\textbf{从任意物点读出共同中心的终像}

两薄透镜的有向法向为 \(\boldsymbol n_1,\boldsymbol n_2\)，焦距为
\(f_1,f_2\)。在共同光心的近轴理想模型内，任意物点 \(\boldsymbol r\)
经两镜后为 \[\boldsymbol r_f=
\frac{\boldsymbol r}{1+\boldsymbol n_1\cdot\boldsymbol r/f_1+
\boldsymbol n_2\cdot\boldsymbol r/f_2}.\]
终像与物点、共同光心共线。像为无穷远的物点满足一条直线
\[1+\left(\frac{\boldsymbol n_1}{f_1}+
\frac{\boldsymbol n_2}{f_2}\right)\cdot\boldsymbol r=0.\] 取原题
\(f_1=f_2=f\)，则有效法向和为 \((3/2,-\sqrt3/2)\)，其长度为
\(\sqrt3\)，所以该无穷远物线到光心的垂距为 \(f/\sqrt3\)。 物点
\(A=(-f,f/\sqrt3)\)
在第一镜自身的焦线，首次像为无穷远，但它不在两镜共同的无穷远物线。两镜总分母为
\(-1\)，故终像是 \(-A=(f,-f/\sqrt3)\)，仍有限。
单镜中间像为无穷远时，先用齐次坐标相乘再归一化，可以避免中途除零。这里的次序独立依赖共同中心、零间隔和原题的理想模型；移动任一镜或改为有间隔的系统，就必须重新加入传播矩阵。

\newpage

\section{第 8 题}\label{ux7b2c-8-ux9898}

用均匀电阻丝制作无限阶指尖陀螺网络。第 \(1\) 阶是边长为 \(l\)
的正三角形框；第 \(2\) 阶在第 \(1\)
阶三边的循环三等分点处内接一个正三角形框；第 \(3\) 阶再在第 \(2\)
阶中按同样规则内接，依次无限进行。电阻丝单位长度质量为
\(\lambda\)，单位长度电阻为 \(r\)。求：

（1）无限阶网络绕共同重心且垂直于三角形平面的轴转动时的转动惯量 \(J\)；

（2）第 \(2\) 阶网络端点 \(A\)、\(B\) 间的等效电阻 \(R_2\)；

（3）无限阶网络端点 \(A\)、\(B\) 间的等效电阻 \(R_\infty\)。

\begin{figure}[H]
\centering
\includegraphics[width=0.566\linewidth,height=\textheight,keepaspectratio,alt={图8-1　三等分点内接生成的前三级正三角形电阻丝网络}]{../figures/question-08.pdf}
\caption*{图8-1　三等分点内接生成的前三级正三角形电阻丝网络}
\end{figure}

\subsection{一、题型定位}\label{ux4e00ux9898ux578bux5b9aux4f4d-7}

三等分接点先决定内三角形缩放比。转动惯量逐层相加，电阻则先把内层换成等效三端网络，再建立递推。

\subsection{二、知识点}\label{ux4e8cux77e5ux8bc6ux70b9-7}

\begin{enumerate}
\def\labelenumi{\arabic{enumi}.}
\tightlist
\item
  三等分点坐标与内接正三角形边长。
\item
  均匀细杆关于平面外轴的转动惯量。
\item
  平行轴定理与几何级数。
\item
  三端对称网络的等效三角形表示。
\item
  节点电势法、递推关系和稳定不动点。
\end{enumerate}

\subsection{三、物理图像}\label{ux4e09ux7269ux7406ux56feux50cf-7}

设外三角形顶点为

\[
A=(0,0),\quad B=(l,0),\quad C=\left(\frac l2,\frac{\sqrt3l}{2}\right). \tag{1}
\]

按图示循环方向取三个三等分点，例如

\[
P=\left(\frac{2l}{3},0\right),\quad
Q=\left(\frac{2l}{3},\frac{\sqrt3l}{3}\right),\quad
R=\left(\frac l6,\frac{\sqrt3l}{6}\right). \tag{2}
\]

计算任一内边：

\[
PQ=\sqrt{0^2+\left(\frac{\sqrt3l}{3}\right)^2}
=\frac l{\sqrt3}. \tag{3}
\]

因此相似缩放比为

\[
q=\frac1{\sqrt3}. \tag{4}
\]

\begin{figure}[H]
\centering
\includegraphics[width=0.809\linewidth,height=\textheight,keepaspectratio,alt={图8-2　循环三等分节点与内层等效三端网络}]{../figures/solution-08.pdf}
\caption*{图8-2　循环三等分节点与内层等效三端网络}
\end{figure}

\subsection{四、解题入口}\label{ux56dbux89e3ux9898ux5165ux53e3-7}

各级三角形具有同一个重心，且第 \(n\) 级新增三角形框边长为

\[
s_n=lq^{n-1}=l\left(\frac1{\sqrt3}\right)^{n-1}. \tag{5}
\]

电阻部分中，任一具有三重旋转对称的线性三端网络，从外部三个端点看，都可等效为三条相同电阻组成的
\(\Delta\) 形网络。记第 \(n\) 阶网络的等效 \(\Delta\) 边电阻为
\(X_n\)，则端点 \(A\)、\(B\) 间的二端等效电阻为

\[
R_n=X_n\parallel(2X_n)=\frac23X_n. \tag{6}
\]

\subsection{五、解析}\label{ux4e94ux89e3ux6790-7}

\needspace{3\baselineskip}
\textbf{（一）单个边长为 \(s\) 的三角形框转动惯量}

每条边质量为 \(\lambda s\)。一条边关于自身中点、垂直平面的轴的转动惯量为

\[
I_{\rm cm}=\frac1{12}(\lambda s)s^2
=\frac{\lambda s^3}{12}. \tag{7}
\]

等边三角形重心到任一边中点的距离等于内切圆半径：

\[
d_c=\frac{\sqrt3}{6}s. \tag{8}
\]

其平方为

\[
d_c^2=\frac{s^2}{12}. \tag{9}
\]

由平行轴定理，一条边关于三角形重心轴的转动惯量为

\[
\begin{aligned}
I_{\rm side}
&=I_{\rm cm}+(\lambda s)d_c^2\\
&=\frac{\lambda s^3}{12}
+\lambda s\cdot\frac{s^2}{12}\\
&=\frac{\lambda s^3}{6}.
\end{aligned} \tag{10}
\]

三条边相加：

\[
\boxed{I_\triangle(s)=3I_{\rm side}=\frac12\lambda s^3}. \tag{11}
\]

\needspace{3\baselineskip}
\textbf{（二）无限阶转动惯量}

各级新增三角形只在有限个节点相交，不存在有长度的重叠线段，因此转动惯量可逐级相加：

\[
\begin{aligned}
J
&=\sum_{n=1}^{\infty}I_\triangle(s_n)\\
&=\frac12\lambda l^3
\sum_{n=1}^{\infty}q^{3(n-1)}.
\end{aligned} \tag{12}
\]

由

\[
q^3=\left(\frac1{\sqrt3}\right)^3
=\frac1{3\sqrt3}, \tag{13}
\]

几何级数求和：

\[
\begin{aligned}
J
&=\frac12\lambda l^3
\frac1{1-1/(3\sqrt3)}\\
&=\frac12\lambda l^3
\frac{3\sqrt3}{3\sqrt3-1}.
\end{aligned} \tag{14}
\]

所以

\[
\boxed{J=\frac{3\sqrt3}{2(3\sqrt3-1)}\lambda l^3}. \tag{15}
\]

\needspace{3\baselineskip}
\textbf{（三）第 \(2\) 阶网络的等效电阻}

先把外三角形每条边总电阻记为

\[
R=rl. \tag{16}
\]

内三角形每条边长 \(l/\sqrt3\)，电阻为

\[
x=\frac{rl}{\sqrt3}=\frac R{\sqrt3}. \tag{17}
\]

外边被连接点分成 \(2R/3\) 与 \(R/3\)
两段。循环三等分网络具有三重旋转对称而没有这里所说的镜面对称。线性无源电阻网络还满足互易性，故三端导纳矩阵等效为等边\(\Delta\)；在A、B间加电压U且C开路时，\(V_C=U/2\)。对其余三个连接节点列基尔霍夫方程并消元，可得从
\(A\) 流入的总电流

\[
I_A=\frac{9(x+R)}{2R(3x+2R)}U. \tag{18}
\]

于是

\[
R_2=\frac{U}{I_A}
=\frac{2R(3x+2R)}{9(x+R)}. \tag{19}
\]

代入 \(x=R/\sqrt3\)：

\[
\begin{aligned}
R_2
&=\frac{2R\left(\sqrt3R+2R\right)}
{9\left(R/\sqrt3+R\right)}\\
&=\frac{2R(\sqrt3+2)}{9(1+1/\sqrt3)}\\
&=\frac{2R\sqrt3(\sqrt3+2)}{9(\sqrt3+1)}.
\end{aligned} \tag{20}
\]

有理化或直接约分可得

\[
\frac{2\sqrt3(\sqrt3+2)}{9(\sqrt3+1)}
=\frac{3+\sqrt3}{9}. \tag{21}
\]

所以

\[
\boxed{R_2=\left(\frac13+\frac{\sqrt3}{9}\right)rl
=\frac{3+\sqrt3}{9}rl}. \tag{22}
\]

\needspace{3\baselineskip}
\textbf{（四）建立无限网络的递推式}

记第 \(n\) 阶三端网络的等效三角形边电阻为 \(X_n\)。第 \(1\)
阶本身就是每边电阻 \(R\) 的三角形，故

\[
X_1=R. \tag{23}
\]

在第 \(n\) 阶中，内层是缩小 \(q=1/\sqrt3\) 倍的第 \(n-1\)
阶网络。电阻与所有长度同比缩放，所以内层等效三角形的边电阻为

\[
x_n=qX_{n-1}=\frac1{\sqrt3}X_{n-1}. \tag{24}
\]

由式（19），第 \(n\) 阶的二端电阻为

\[
R_n=\frac{2R(3x_n+2R)}{9(x_n+R)}. \tag{25}
\]

又因 \(R_n=2X_n/3\)，所以

\[
\begin{aligned}
X_n
&=\frac32R_n\\
&=\frac{R(3x_n+2R)}{3(x_n+R)}.
\end{aligned} \tag{26}
\]

引入无量纲量

\[
y_n=\frac{X_n}{R}. \tag{27}
\]

由 \(x_n/R=y_{n-1}/\sqrt3\)，式（26）化为

\[
\boxed{
y_n=\frac{\sqrt3 y_{n-1}+2}
{\sqrt3 y_{n-1}+3}}\qquad y_1=1. \tag{28}
\]

\needspace{3\baselineskip}
\textbf{（五）求递推的不动点}

当 \(n\to\infty\) 时，递推收敛到正不动点 \(y\)：

\[
y=\frac{\sqrt3y+2}{\sqrt3y+3}. \tag{29}
\]

交叉相乘：

\[
y(\sqrt3y+3)=\sqrt3y+2. \tag{30}
\]

展开并移项：

\[
\sqrt3y^2+(3-\sqrt3)y-2=0. \tag{31}
\]

由求根公式：

\[
y=\frac{-(3-\sqrt3)\pm
\sqrt{(3-\sqrt3)^2+8\sqrt3}}
{2\sqrt3}. \tag{32}
\]

根号内逐步化简：

\[
\begin{aligned}
(3-\sqrt3)^2+8\sqrt3
&=9-6\sqrt3+3+8\sqrt3\\
&=12+2\sqrt3.
\end{aligned} \tag{33}
\]

电阻必须为正，舍去负根：

\[
y=\frac{\sqrt3-3+\sqrt{12+2\sqrt3}}{2\sqrt3}. \tag{34}
\]

最后

\[
R_\infty=\frac23X_\infty
=\frac23Ry. \tag{35}
\]

代入 \(R=rl\) 与式（34）：

\[
\boxed{
R_\infty
=\frac{\sqrt3-3+\sqrt{12+2\sqrt3}}
{3\sqrt3}\,rl}. \tag{36}
\]

其数值系数约为 \(0.51278\)，小于第 \(1\) 阶的 \(2/3\)，也略小于第 \(2\)
阶的
\((3+\sqrt3)/9\simeq0.52578\)，符合不断加入并联通路后等效电阻下降的直觉。

式（18）可通过节点电势法完整核查。令 \(V_A=U\)、\(V_B=0\)，未知节点取
\(V_C,V_P,V_Q,V_R\)。由三端旋转对称与互易性得 \(V_C=U/2\)；其余方程解为

\[
V_P=V_Q=\frac{x+R}{3x+2R}U,\qquad
V_R=\frac{5x+2R}{6x+4R}U,
\]

再把与 \(A\) 相连支路电流相加即可得到式（18）。这里把内层网络换成等边
\(\Delta\)
不是近似，而是线性、无源、三端且具有三重旋转对称时的严格端口等效。

递推收敛性也可核验。设 \(F(y)=(\sqrt3y+2)/(\sqrt3y+3)\)，则

\[
F'(y)=\frac{\sqrt3}{(\sqrt3y+3)^2}.
\]

在正数范围内 \(0<F'(y)<\sqrt3/9<1\)，因此正不动点具有吸引性，从
\(y_1=1\) 出发必收敛。

单个三角形框的式（11）也可由连续积分复核。以一条水平边为例，边中点到三角形重心的距离为
\(d_c\)，沿边取坐标 \(\xi\in[-s/2,s/2]\)，质量元为
\(\mathrm dm=\lambda\,\mathrm d\xi\)。该边贡献

\[
\begin{aligned}
I_{\rm side}
&=\int_{-s/2}^{s/2}(\xi^2+d_c^2)\lambda\,\mathrm d\xi\\
&=\lambda\left[\frac{\xi^3}{3}\right]_{-s/2}^{s/2}
+\lambda d_c^2s\\
&=\frac{\lambda s^3}{12}+\frac{\lambda s^3}{12}\\
&=\frac{\lambda s^3}{6}.
\end{aligned}
\]

它与平行轴定理完全一致，也说明为何转动惯量缩放为长度的三次方。

对递推式作数值前两步核验：\(y_1=1\)，

\[
y_2=\frac{\sqrt3+2}{\sqrt3+3}
=\frac{3+\sqrt3}{6}\simeq0.788675,
\]

所以 \(R_2/R=(2/3)y_2\simeq0.525783\)。再迭代一次得
\(R_3/R\simeq0.51468\)，已接近
\(R_\infty/R\simeq0.51278\)。这组数值既显示收敛很快，也可及时发现漏乘
\(2/3\) 的常见错误。

若需要从头复算式（18），可把六个外边分段依次记为 \(2R/3,R/3\)，并对
\(P,Q,R\) 三个连接点列方程。例如 \(P\) 与 \(A\)、\(B\)
及两个内层节点相连，故

\[
\frac{V_P-U}{2R/3}
+\frac{V_P-0}{R/3}
+\frac{V_P-V_Q}{x}
+\frac{V_P-V_R}{x}=0.
\]

\(Q\)、\(R\) 方程由循环置换得到，\(C\)
点方程再联立。解出各节点电势后，从 \(A\) 流入的电流是与 \(A\)
相连两支路电流之和：

\[
I_A=\frac{U-V_P}{2R/3}
+\frac{U-V_R}{R/3}.
\]

代入节点解即可化为式（18）。列式时每个外顶点只连两段外边，内层节点之间才连电阻
\(x\)；若误把内层等效三角形直接并在 \(AB\)
两端，就会丢失三等分接点的拓扑信息。

无限网络的答案还满足 Rayleigh
单调性：在既有电阻网络上增加有限正电阻支路，不会增大任意两端的等效电阻。因此

\[
R_1\ge R_2\ge R_3\ge\cdots>0.
\]

递推函数 \(F(y)\) 单调递增，且对 \(y\) 大于正不动点时有 \(F(y)<y\)；从
\(y_1=1\) 出发，序列严格下降并以下界 \(y_\infty\)
收敛。这样既从电路定理，又从代数递推两条路径确认应选择二次方程的正根。若计算结果出现
\(R_\infty>R_1\) 或负值，无需检查复杂化简便可立即判错。

转动惯量也有上下界检查。只保留第一级时
\(J_1=\lambda l^3/2\)；所有更小三角形提供正贡献，故 \(J>J_1\)。同时公比
\(1/(3\sqrt3)<1/5\)，所以

\[
J<\frac{J_1}{1-1/5}=\frac58\lambda l^3.
\]

闭式系数 \(3\sqrt3/[2(3\sqrt3-1)]\simeq0.6192\) 的确位于 \(0.5\) 与
\(0.625\) 之间。

\needspace{3\baselineskip}
\textbf{端口等效要检查功率，而不只检查一个电阻}

取第2阶\(R=1\ \Omega\)、内三角形边\(x=1/\sqrt3\ \Omega\)、端口电压\(U=1\ {\rm V}\)。节点解给出
\[V_C=\frac12,\quad V_P=V_Q=1-\frac1{\sqrt3},\quad
V_R=-\frac12+\frac2{\sqrt3}.\]
电势值均位于0和1之间，满足无源电阻网络的最大值性质。P、Q等势使PQ支路电流为零，但其余内边仍有电流，不能据此删掉整个内三角形。

由A接出的两支路电流为 \[I_{AP}=\frac{1-(1-1/\sqrt3)}{2/3},\qquad
I_{AR}=\frac{1-(-1/2+2/\sqrt3)}{1/3}.\] 两者相加得到
\[I_A=\frac9{3+\sqrt3}\ {\rm A},\]
恰为\(1/R_2\)。所有九条电阻支路的耗散总功率也应满足
\[\sum_{\langle ij\rangle}\frac{(V_i-V_j)^2}{R_{ij}}
=UI_A.\]
这个恒等式由节点方程相乘并求和得到，每条无向边只数一次。它既检验三等分电阻的接法，又检验端点电流没有漏支路。图8-2保留了循环三等分方向，不能为了套用镜面对称把接点改成另一组。

\subsection{六、关键洞察}\label{ux516dux5173ux952eux6d1eux5bdf-7}

相似缩放对不同物理量的幂次不同：质量随长度为 \(q\)，到重心轴的距离平方随
\(q^2\)，所以转动惯量随 \(q^3\)；电阻只随长度为
\(q\)。同一几何递归因此在两问中产生不同的级数或迭代。

\subsection{七、易错点}\label{ux4e03ux6613ux9519ux70b9-7}

\begin{enumerate}
\def\labelenumi{\arabic{enumi}.}
\tightlist
\item
  把内三角形边长直接当成 \(l/3\)。
\item
  单边转动惯量只用 \(Ms^2/12\)，漏掉平行轴项。
\item
  把各级网络当成简单并联，忽略内三角形接在外边三等分节点上。
\item
  将三端网络的 \(\Delta\) 边电阻 \(X_n\) 与两端等效电阻 \(R_n=2X_n/3\)
  混淆。
\item
  求二次方程后未舍去负电阻根，也未用单调性核查结果。
\end{enumerate}

\subsection{八、课堂讲解建议}\label{ux516bux8bfeux5802ux8bb2ux89e3ux5efaux8bae-7}

先用坐标计算一次内边长度，避免学生从图形比例猜
\(1/3\)。电阻部分先把``任意对称三端网络''画成等边三角形，只讨论 \(X_n\)
与 \(R_n\) 的区别；随后再把内层缩放接回外层。最后用
\(R_1>R_2>R_\infty>0\) 做整体检查。

\subsection{九、变式训练}\label{ux4e5dux53d8ux5f0fux8badux7ec3-7}

若每一级改在边中点内接倒三角形，缩放比变为
\(1/2\)，转动惯量级数的公比变为
\(1/8\)；电阻网络的连接拓扑也随之改变。可要求学生只建立递推而不必求闭式，以训练``几何缩放---端口等效---不动点''的通用方法。

\needspace{3\baselineskip}
\textbf{截断阶数需要多少}

转动惯量的第n阶截断误差可精确写为 \[J_\infty-J_n
=\frac{\lambda l^3}{2}\frac{q^{3n}}{1-q^3},\qquad q=1/\sqrt3.\]
所以相对总惯量误差就是\(q^{3n}\)。n=3时为\(1/(81\sqrt3)\simeq0.007128\)，三层已把惯量误差降到1\%以下。

电阻递推的统一收缩上界是\(k=\sqrt3/9\)：
\[|y_n-y_\infty|\le k^{n-1}|1-y_\infty|.\]
对二端电阻同乘\(2R/3\)即可。这个上界比实际局部收敛率保守，却可在不知道后续节点电势时使用。惯量误差由几何级数精确给出；电阻误差由非线性递推界给出，两者不能共用同一个\(q^3\)。

若比较数值，第3阶二端电阻约\(0.51468R\)，无限阶约\(0.51278R\)，剩余误差小于第2阶到无限阶的误差，符合新增正支路的单调性。

\needspace{3\baselineskip}
\textbf{逐层电阻与两端电阻的误差尺度}

令 \(\Delta\) 等效边电阻为 \(X_n=Ry_n\)。内层边长缩为
\(1/\sqrt3\)，所以传给外层的边电阻是
\(x_n=Ry_{n-1}/\sqrt3\)。代入节点法所得两端电阻并乘 \(3/2\)：
\[y_n=\frac{\sqrt3y_{n-1}+2}{\sqrt3y_{n-1}+3}=F(y_{n-1}).\] 其导数为
\[F'(y)=\frac{\sqrt3}{(\sqrt3y+3)^2}>0.\] 因此误差不交替变号。从
\(y_1=1\) 开始的各阶都在正不动点上方，逐层下降。若把电路两端电阻
\(R_n=2Ry_n/3\) 直接当作内层三角形的边电阻，会漏掉 \(3/2\) 因子。

用不动点 \(y_*\) 写出精确差： \[F(y)-F(y_*)=\frac{\sqrt3(y-y_*)}
{(\sqrt3y+3)(\sqrt3y_*+3)}.\] 所以每层的实际误差比例随 \(y\)
略变，并不是固定的等比级数。在 \(y\ge0\) 上可用统一上界
\[|y_n-y_*|\le\left(\frac{\sqrt3}{9}\right)^{n-1}|y_1-y_*|.\]
转动惯量因各层质量和长度平方直接相乘，其尾项是严格几何级数；电阻来自节点耦合，只能从递推映射导出误差比例。两种收敛机制不可混用。
\end{document}
